%-----------------------------------------------------------------------------
% Template for AIMS Structured Masters Research Project
%
% The fonts, linespacing, numbering, page styles, order
% of  Title/Abstract/TOC/Body/{Appendices}/Acknowledgements/References 
% are prescribed as the AIMS house style.
%
% Do not change them or add to it without getting approval first.
% Essays are not accepted if they do not follow house style.
% This is in preparation for your Masters where the university
% will be much more strict on the house style.
%
\documentclass{aimsessay}
% \usepackage{etex} % fix for using xy and tikz in the same document
\usepackage[utf8]{inputenc}
\usepackage[T1]{fontenc}
\usepackage{lmodern}
\usepackage[round]{natbib}
\usepackage{amsmath}
\usepackage{amsfonts}
\usepackage{amssymb}
\usepackage{mathtools}
\usepackage{latexsym}
\usepackage{parskip}
\usepackage{enumerate}
\usepackage{tikz}
\usepackage{lipsum}
%\DeclareUnicodeCharacter{U+2032}
% Fix spacing before theorem environments
\makeatletter\def\thm@space@setup{%
  \thm@preskip=1.0\parskip \thm@postskip=0pt
  }
%\makeatother
% Reduce space before proof
%\makeatletter
\renewenvironment{proof}[1][\proofname]{\par
  \vspace{-\topsep}% remove the space after the theorem
  \pushQED{\qed}%
  \normalfont
  \topsep0pt \partopsep0pt % no space before
  \trivlist
  \item[\hskip\labelsep
        \itshape
    #1\@addpunct{.}]\ignorespaces
}{%
  \popQED\endtrivlist\@endpefalse
  \addvspace{0pt plus 0pt} % some space after
}
%\makeatother
% Reduce spacing around section headings?
%
%-----------------------------------------------------------------------------
% To use external packages for specific needs, 
% get approval before adding them here (they
% should not override general AIMS house style and layout).
%
% Examples:
% 
% Load caption before arabtex
\usepackage{caption} %many figures in one figure (note subfigure and subfig are deprecated) 
\usepackage{subcaption} %many figures in one figure (note subfigure and subfig are deprecated) 
% For Arabic
%\usepackage{arabtex}
%\usepackage{utf8}
%\setcode{utf8}
% For tables:
\usepackage{booktabs} % \toprule, \midrule, \bottomrule
\usepackage{array}    % \newcolumntype
% 
% For figures:
%\usepackage[below,section]{placeins} % use \FloatBarrier in the body to really force a float somewhere. Please limit use.
\usepackage{float}  %"H" placement spec, for **really here** (i.e. not float)
%
% For code and algorithms
\usepackage{moreverb}   % \verbatimtabinput
% \usepackage{listings} % more flexible and complicated 
                        % than moreverb and algorithm
%\usepackage[maxlevel=4]{csquotes}
% 
% Others
% \usepackage[all]{xy}  %  if you use this uncomment \usepackage{etex} above to fix a conflict with tikz.
% \usepackage{sagetex}
% \usepackage{siunitx} % to typeset numbers, units, align decimals in tables.
% \usepackage{dcolumn} % less flexible but maybe faster than siunitx above.
% \usepackage{mathtools} % More maths, e.g. \mathclap.
%
% Others may be landscape, longtable, algorithm, algorithmic, etc.
% 
% ----------------------------------------------------------------------------
% An AIMS Essay can use the sectioning commands "\chapter", "\section",
% "\subsection". No "\subsubsection", "\paragraph", etc. They are disabled.
% 
% For Theorems and such, use the environments defined here:
% \begin{thm}...\end{thm} (or "lem", "defn", etc)
% 
% We put the number to the left of the Theorem heading.
%\swapnumbers 
% 
% Theorems are in italics.
\theoremstyle{plain}
\newtheorem{thm}[subsection]{Theorem}
%
% Rest is not in italics.
\theoremstyle{definition} 
\newtheorem{lem}[subsection]{Lemma}
\newtheorem{cor}[subsection]{Corollary}
\newtheorem{conj}[subsection]{Conjecture}
\newtheorem{pro}[subsection]{Proposition}
\newtheorem{exa}[subsection]{Example}
\newtheorem{defn}[subsection]{Definition}
\newtheorem{rem}[subsection]{Remark}
\newtheorem{Rem}[subsection]{Conclusions}
\DeclareMathOperator{\sech}{sech}


%\DeclareRobustCommand{\rchi}{{\mathpalette\irchi\relax}}
% 
% If you have no theorems, but a lot of equations, maybe the
% following two lines are good. Beware of corresponding Equation
% and Example numbers though! Number equations by sections.
% 
\numberwithin{equation}{section}
%
%----------------------------------------------------------------------------- 
% Abstracts are usually written in English, with a version in your
% mother tongue underneath.
%
% French, Igbo, Malagasy, etc. students use normal LaTeX
% for special characters, for example \'{e}
%
% Amharic students use LibreOffice to write Amharic,
% export and include a figure.
%
%\begin{RLtext}    
%Here is where the arabic text goes.
%You can just type it with an arabic keyboard
%\end{RLtext}\\
%-----------------------------------------------------------------------------
 
% Your own command shortcuts can go here
% keep them clearly separate from other sections of the preamble
% It is good style to have only a few of these so that
% we can read one another's code. If you have to many, 
% then your code does not compile in someone else's template easily,
% and it makes it harder to read. These definitions are only
% meant for very often-used commands to save typing; Examples:
%
%\newcommand {\be}{\begin{equation}}
%\newcommand {\ee}{\end{equation}}
%\newcommand {\C}{\mathbb{C}} % Complex
%\newcommand {\Z}{\mathbb{Z}} % Integers
%\newcommand {\R}{\mathbb{R}} % Real
%\DeclareMathOperator{\sech}{sech} % declaring new math operators like \sin.
%  
%-----------------------------------------------------------------------------
% Title & Author
% On this page you must have NO full-word capitalizations, bold, or colour.  
% All AIMS research projects per year have the same title page.
% In English your family name is written last, i.e. Firstname LASTNAME
% English Capitalization, not as in some Francophone countries where
% you write LASTNAME, Firstname.
% Put your AIMS email address only please, for consistency,
% not gmail or some other webmail address.
\title{Lie Group Analysis of a Nonlinear  Coupled System of \newline Korteweg-de Vries Equations}
% Your name must be in English Capitalisation with no comma, 
% and the Family name comes last.
\author{Joseph Owuor Owino(josephowino@aims.ac.za)\\
% Then in the MAIN BODY use this:                  
%\begin{RLtext}
%نووووووسسسسح
%\end{RLtext}\\
%%%%%%%%%%%%%%%%%%%%%%%%%%%%%%%%%%% 
%\begin{otherlanguage}{arabtext}
%شةشىغ
%\end{otherlanguage}\\
%%%%%%%%%%%%%%%%%%%%%%%%%%%%%%%%%%%%
% Amharic students speak to me about how to add your name in your own alphabet.
% Everything here is prescribed; do not enter bold or ALL CAPS here,
% it will not be accepted.
African Institute for Mathematical Sciences (AIMS)\\
\\
% Example1
% Supervised by Professor Barry Green
% University of Stellenbosch, South Africa
% Example2
% Supervised by Doctor Ingrid Rewitzky
% University of Stellenbosch, South Africa
{\small Supervised by: Professor Chaudry Masood Khalique}\\
{\small North-West University, South Africa}%
% Don't put the department, it becomes too long.
}
\date{{\small 14 May  2020}\\%
  {\scriptsize\it Submitted in partial fulfillment of 
    a structured masters degree at AIMS South Africa}\\%
  \vspace{0.5cm}{\includegraphics{images/AIMS_SA_Logo.pdf}}}
%-------------------------------------------------------------------------
\begin{document}
%\selectlanguage{english}
\pagestyle{empty}
\maketitle
% All other files are included via \input. 
% To compile in texmaker while viewing any of those
% without having to switch back, use
%   Options > Define Current Document as 'Master Document'
% To not have to close a PDF, remove viewpdf from quickbuild 
% and open the PDF (once) manually: it will auto-refresh or with control-r
% 
%-------------------------------------------------------------------------
% The abstract is the first thing we want to see. No acknowledgements or 
% dedications here. Fetch the abstract from a separate file.
% Please write it in English and in your mother tongue.
% An abstract should be less than half a page, so that both abstracts 
% (that is both languages) fit onto one page.
% We number roman numerals until the main body
\pagenumbering{roman}
% Abstracts are usually written in English, with a version in your
% mother tongue underneath
\chapter*{Abstract} 
\addcontentsline{toc}{chapter}{Abstract}
% Don't change anything above this.
Lie group analysis of nonlinear partial differential equations (NLPDEs) makes use of symmetries of the NLPDE to find solutions and conservation laws.  The solutions describe the state of the system at all times and conservation laws give its invariant properties. NLPDEs have applications in many settings. However, finding exact solutions to such NLPDEs is a challenge because of the increased complexity since NLPDEs have many independent variables. In this dissertation, we propose to research on  Lie group analysis of a  nonlinear coupled system of Korteweg de-Vries (KdV) equations. To achieve this, we first introduce the mathematical concepts of Lie group analysis and conservation laws. As an illustrative example, we perform a Lie group analysis of the KdV equation. This means we find Lie point symmetries and compute invariant solutions. Conservation laws for mass, momentum, and energy are also constructed for the KdV equation by the method of multipliers. Next, we focus on Lie group analysis of a nonlinear system of coupled KdV equations. We find symmetries of the coupled system and use them for symmetry reductions. We also compute invariant solutions for the coupled KdV equations. Furthermore, conservation laws are constructed by using multipliers and a theorem proposed by Nail Ibragimov. The target application for this dissertation is in systems that can be described by a KdV model. Particularly systems behaving like solitons which are traveling waves that do not change shape after a pair-wise collision.


\textbf{Keywords:} coupled KdV equations; Lie group analysis; symmetry reductions; soliton; multipliers; \newline conservation laws.  
%You may like to repeat the abstract icoupled KdV equations; Lie group analysis, symmetry reductions; travelling wave solutions; multipliers, conservation laws. n your mother tongue.

% At a unviersity like Stellenbosch you *must* produce an abstract in Afrikaans for your masters.
% At AIMS you are encouraged to repeat the abstract in your mother tongue
% French, Igbo, Mlagasy, etc. just write it using LaTeX's special
% characters.
% Arabic students see the arabic.tex file for an example
% Amharic use openoffice and export from there and import a figure here.
% Where the words do not exist put the English work in italics, or use mathematical symbols.


% Do not change anything below this except for adding your
% signature (replace images/signature.png) and your name.
\vfill
\section*{Declaration}
I, the undersigned, hereby declare that the work contained in this research project is my original work, and that any work done by others or by myself previously has been acknowledged and referenced accordingly.

% Scan your signature into a small picture called 'signature.png' and insert it
% above your name and the date:
\includegraphics[height=2cm]{images/signature.png} \newline \hrule
% Your name must be in English Capitalisation with no comma, and the Family name comes last. 
% Do note the date below. It is called the "deadline".
Joseph Owuor Owino, 14 May 2020


% Don't go typing out the contents.
\tableofcontents
%\listoftables 
\newpage
% We switch to arabic numerals here where your page count starts
% Essays must be close to 30 pages long *starting here* and up to and including
% the conclcusion. It does not include the acknowledgements or references.
% 
% Figures may differ between topics, but they are not there
% to fill the pages -- they must add meaning.
% In general most figures should be 0.8 times the width of the page
% (perhaps wider in total when two or three columns of figures)
% See the example in chapter one for defining that. Be *consistent*
% in your presentation of information.
\pagenumbering{arabic}
\pagestyle{myheadings}
%-----------------------------------------------------------------------------
% Each chapter goes in a separate file
% Chapter titles are just examples
% Always have a question
% Note the Case Pattern used at AIMS
<<<<<<< HEAD
\chapter{Introduction \label{ch1}}
Nonlinear partial differential equations (NLPDEs) are used in the formulation of many natural laws such as in fluid dynamics. NLPDEs have been very key to technological developments as we look for numerical approximations and simulations using computers. To be able to this, we must be able to obtain approximations to solutions or exact closed-form solutions. In order to understand NLPDEs and the systems they represent better, several methods have been proposed and are still being used to describe the state of NLPDE systems. Lie group analysis is one of those methods.

Lie group analysis was inspired was the work Sophus Lie \citep{lie1891vorlesungen}  about applying Galois theory of algebraic equations to differential equations (DEs). He showed that Lie group-invariant DEs could be reduced to a simpler form that is easily solvable. Furthermore, group-invariant solutions of the equation could also be found. There are many books written on this subject such as  \citep{ovsyannikov2013lectures}, \citep{olver2000applications}, \citep{bluman2013symmetries}, \citep{ibragimov}, \citep{ibragimov2009practical} and the method has been recently extensively used in solving NLPDEs.

Conservation laws, by which we mean a divergence expression vanishing on all the solutions of a system and describes invariant properties of the system, are very important in the understanding of real physical phenomena. They have applications in various mathematical sciences disciplines. Conservation laws have been used to explain the integrability of DEs and the effectiveness of numerical methods \citep{leveque1992numerical}. Recently,  conservation laws have been used to construct exact closed-form solutions of certain partial differential equations (PDEs). This needs a good understanding of conservation laws.

Emmy Noether  \citep{noether1918konig} found a connection between conservation laws and symmetries for DEs that have a Lagrangian which is a function that describes the dynamics of the system. For those DEs that do not have a Lagrangian, methods which do not rely on the variational principle have been developed. Some of these methods include the multipliers approach and a conservation theorem proposed by Nail Ibragimov \citep{ibragimov2007new}.

To contribute to Lie group analysis of NLPDEs, we shall in this dissertation investigate an application to a  nonlinear coupled system of Korteweg-de Vries (KdV) equations. The KdV equation describes the dynamics of solitons, and ion-acoustic waves in plamsas. We shall perform a Lie group analysis of the system, that is, finding Lie point symmetries as well as group-invariant solutions and conservation laws.

The application addressed in this dissertation is an understanding of a typical KdV equation which has uses in many areas of mathematical sciences. Some of the many applications include describing the dynamics of shallow-water waves,  ion-acoustic waves in plasmas, and long internal waves in oceans.

To reach our goal of contributing to Lie group analysis of PDEs, we shall in Chapter \ref{ch2} develop mathematical concepts of Lie group analysis and present some algorithms for computing conservation laws. In Chapter \ref{ch3}, we will present an illustrative example of Lie group analysis of the KdV equation. That means we find Lie point symmetries, group-invariant solutions including a one-soliton solution, and compute conservation laws using multipliers. Next, we will in Chapter \ref{ch4} investigate a system of coupled KdV equations. We find Lie point symmetries of the coupled KdV system and use them for symmetry reductions. We also compute invariant solutions. Moreover, we find conservation laws for the nonlinear coupled system of  KdV equations by using multipliers and a theorem proposed by Nail Ibragimov. We conclude the work in Chapter \ref{ch5} and suggest further work.
=======
\chapter{Prelude}
Has it ever occurred to you that there exists some underlying theory that unifies a bewildering  variety of special ad-hoc techniques designed by researchers to solve certain particular seemingly, unrelated types of differential equations such as separable, homogeneous or exact equations? Well, this seems to have been  the state of the art around the middle of nineteenth century \citep{olver2000applications},\citep{arrigo2015symmetry}. 


Incited by the impetus to extend Evariste Galois's theory on algebraic equations to differential equations, a heavy Norwegian mathematician, Marius Sophus Lie(1832-1899) \citep{lie1891vorlesungen},\citep{fritzsche1999sophus}  initiated a monumental study of continuous transformation groups, presently known as Lie groups. Lie made a profound and far reaching discovery that these special methods, were in fact, all special cases of a general integration procedure based on invariance of differential equations under a continuous group of symmetries \citep{olver2000applications}.
% cite Galois book and Sophus lie paper

Roughly speaking, a symmetry group of a system of partial differential equations is a group which transforms solutions of the system to another solution of the same system. From  a classical framework, these groups consist of geometric transformations on the space of independent and dependent variables, and act on solutions, by transforming their graphs. Typical examples are groups of translations, rotations and reflections.\citep{olver2000applications}
%%% Insert image of showing rotation , reflection and translation
Lie groups have had a profound impact on all areas of mathematics, both pure and applied, as well as physics, engineering and other mathematically-based sciences. The applications of Lie's continuous symmetry groups include such diverse fields as algebraic topology, differential geometry, invariant theory, birfucation theory, special functions, numerical analysis, control theory, classical mechanics, quantum mechanics, relativity, continuum mechanics, fluid mechanics, gas dynamics, combustion theory and so on. One can never overestimate Lie's contribution to modern science and mathematics.\citep{olver2000applications}

For the case of ordinary differential equations, invariance under a one-parameter symmetry group implies that we can reduce the order of the equation by one, thereby recovering solutions to the original equations from those of the reduced equation by a single quadrature. This idea has been extended to reducing non-linear partial differential equations invariant under a Lie symmetry to a simpler form with fewer variables and a lower order, hence easy to solve. Precisely speaking, the reduction of a partial differential equation with respect to $\jmath$-dimensional (solvable) subalgebra of its Lie symmetry leads to reducing the number of independent variables by $\jmath$. In addition, similarity or invariant solutions of the partial differential equation can also be established.

Fundamental natural laws of real world and technological problems whose formulations are non-linear ordinary and partial differential equations can be successfully treated and solved by Lie group methods, even when other methods of integration  fail the test. In fact, group analysis is the only universal and effective method for solving non-linear differential and partial equations analytically, dealing equally easily with linear and non-linear equations, as well as with constant coefficient and variable coefficients. It is therefore imperative to study Lie group analysis methods.


For example, from the traditional point of view, the linear equation \citep{ibragimov2009practical} \begin{equation}
\frac{ \mathrm{d}^n y}{ \mathrm{d} x^n } + a_1\frac{ \mathrm{d}^{n-1} y}{ \mathrm{d} x^{n-1} } + \ldots + a_{n-1}\frac{ \mathrm{d} y}{ \mathrm{d} x }+ a_n y =0 \label{1.0.1}
\end{equation} with constant coefficients $ a_1, \ldots,a_n$ is different from the equation \begin{equation}
\bar{x}^n	\frac{ \mathrm{d}^n \bar{y}}{ \mathrm{d} \bar{x}^n } + a_1 \bar{x}^{n-1}\frac{ \mathrm{d}^{n-1} \bar{y}}{ \mathrm{d} \bar{x}^{n-1} } + \ldots + a_{n-1}\bar{x}\frac{ \mathrm{d} \bar{y}}{ \mathrm{d} \bar{x} }+ a_n \bar{y} =0  \label{1.0.2}
\end{equation} known as \textit{Euler} equation. From lie group stand point, however, the equations (\ref{1.0.1}) and (\ref{1.0.2}) are merely two different representations of the same equation with two known commuting symmetries, namely, 
\begin{equation}
\chi_1= \frac{\partial }{ \partial x}, \quad \chi_2= \frac{\partial }{ \partial y} \quad  \text{and} \quad \bar{\chi}_1= \frac{\partial }{ \partial \bar{x}}, \quad  \bar{\chi}_2= \frac{\partial }{ \partial \bar{y}} \label{1.0.3}
\end{equation} for equations (\ref{1.0.1}) and (\ref{1.0.2}) respectively. These two symmetries span two similar Lie algebras and readily lead to the transformation $ x = \ln{\bar{|x|}}$, converting Euler equations to the equation with constant coefficients.
 
 While engrossed in the noblest of a task to find closed form solutions of partial differential equations, one is usually interested in whether or not some quantities remain invariant in the evolution of the partial differential equation. The idea of a \textit{conservation law} which describes quantities that do not dissipate with time. Examples of such conserved quantities in physical systems  include momentum, mass, energy, angular momentum and charge. Most recent use of conservation laws  has been in investigation of existence, uniqueness and stability of non-linear partial differential equations. \citep{lax1968integrals} \citep{benjamin1972stability}
 
 
 Conservation laws have also been applied in development and use of numerical methods \citep{eveque1992numerical},\citep{godlewski2013numerical}. The aforementioned applications emphasize the necessity to study conservation laws. To this end, there is an established intimate connection between symmetries and conservation laws. Thanks to impeccable works of Sophus Lie which inspired a strong German mathematician Ammie Emmy Noether(1882-1935) \citep{tent2008emmy}.
 
 
This connection was unearthed by Emmy Noether and is enshrined in Noether's theorem, "Every differentiable symmetry of the action of a physical system has a corresponding conservation law." This theorem shows how each one parameter variational symmetry group gives rise to a conservation law of the Euler-Lagrange equations. For example, conservation of energy stems from invariance of the problem under a group of time translations while conservation of linear and angular momenta reflect translational and rotational invariance of the system. Noether symmetries have profound applications in the study of properties of particles moving under the influence of gravitational fields \citep{olver2000applications}.

Having noted that some partial differential equations lack a Lagrangian, mathematicians have  developed symmetry  methods of deriving conservation laws that do not rely on the existence of variational principle. Some of these include the multiplier approach \cite{naz2012conservation}, \citep{bluman2010applications}
and another method due to Ibragimov's theorem on adjoint equations of a system of partial differential equations \citep{ibragimov2009practical}.

This dissertation strives to study,....

Lie method is used to reduce .... and construct group invariant solutions of the system. We also derive conservation laws for the system. 

This system of partial differential equations has applications in ... and finding its solutions will improve understanding of physical world processes. The solutions and conservation laws for the case of ... are derived from symmetry point of view. These results will be very important to the mathematical sciences community and the world at large since...

The dissertation  is structured as follows:

Chapter 2 provides a brief and succinct exposition of some concepts on Lie symmetry methods and conservation laws.

Chapter 3 contains an illustration of symmetry methods in action.

Chapter 4 accommodates ....system of equations.

Chapter 5 highlights a summary of work done in this dissertation with mentions of future work.

Bibliography is given at the end.
% Include photos of Soplus Lie and Emmy Noether .

% Reference to Peter J Olver

>>>>>>> f3e458080f41e5290173f14b8121103bb6a456e0
 % Introduction is usually a chapter itself.
<<<<<<< HEAD
\chapter{Preliminaries \label{ch2}}
This chapter serves as a brief introduction of fundamental concepts on Lie group analysis \citep{ibragimov} and conservation laws  \cite{ibragimov2007new}. The first section of the chapter reviews the notions of local Lie groups and infinitesimal transformations, prolongations of Lie groups, group invariants, and symmetry groups. In the second section, we provide an introduction to conservation laws as well as the method of multipliers and Ibragimov's theorem for computing conservation laws. The notations fixed here will be used in the rest of the dissertation.
\section{Lie group analysis} 
This section introduces the basic concepts of the Lie group theory of PDEs and how one can obtain symmetries.
\subsection{Local Lie groups}
The setting is in Euclidean spaces $\mathbb{R}^n$ of $x=x^i$ independent variables and  $\mathbb{R}^m$ of $ u=u^{\alpha}$  dependent variables. Consider the transformations
\begin{align}
T_{\epsilon}: \quad \bar{x}^i = \varphi^i(x^i,u^{\alpha},\epsilon ), \quad \bar{u}^{\alpha } =\psi^{\alpha}(x^i,u^{\alpha},\epsilon), \label{tranfmtn}
\end{align} involving  the continuous  parameter $\epsilon$ which ranges from a neighbourhood $ \mathcal{N}' \subset \mathcal{N} \subset \mathbb{R}$ of $ \epsilon=0 $ where  the functions $ \varphi^i$ and $ \psi^{\alpha}$ differentiable and analytic in the parameter $\epsilon$.
\begin{defn} The set $\mathcal{G}$  of transformations given by  (\ref{tranfmtn}) is a local Lie group if it holds true that
	\begin{enumerate}[(i)]
		\item (Closure) Given \label{propty1} $T_{\epsilon_1},T_{\epsilon_2}  \in \mathcal{G}$, for $\epsilon_1,\epsilon_2 \in  \mathcal{N}' \subset \mathcal{N} $, then $T_{\epsilon_1}T_{\epsilon_2} = T_{\epsilon_3} \in \mathcal{G},\,\,  \epsilon_3= \phi( \epsilon_1,\epsilon_2) \in \mathcal{N}.$	
		\item (Identity) There exists a unique  $  T_0 \in \mathcal{G} $ if and only if $ \epsilon=0$ such that $ T_{\epsilon}T_{0}=T_{0}T_{\epsilon}=T_{\epsilon}$.
		\item (Inverse) There  exists a unique $T_{\epsilon^{-1}} \in \mathcal{G}$ for every transformation $ T_{\epsilon} \in \mathcal{G}, \quad \newline \text{where} \quad \epsilon \in  
		\mathcal{N}' \subset  \mathcal{N}
		 \quad \text{and} \quad  \epsilon^{-1} \in \mathcal{N}
		 $ such that $ T_{\epsilon} T_{\epsilon^{-1}}=T_{\epsilon^{-1}}T_{\epsilon}=T_{0}$.
	\end{enumerate}
\begin{rem}
Associativity property of the group $\mathcal{G} $  defined in (\ref{tranfmtn}) follows from (\ref{propty1}).
\end{rem}
\end{defn}
\subsection{Prolongations\label{pr}}
Since we shall be dealing with PDEs, we wish to describe how the dependant variables $u^{\alpha}$ transform. Consider the  system PDEs
\begin{equation}
\Delta_{\alpha}\left(x^i, u^{\alpha} , u_{(1)}, \ldots,u_{(\pi)}\right) = \Delta_{\alpha}=0. \label{pi-thorderpde}
\end{equation} The partial derivatives $ u_{(1)}=\{ u_{i}^{\alpha}\}, u_{(2)}=\{ u_{ij}^{\alpha}\} , \ldots, u_{(\pi)}=\{ u_{i_1\ldots i_{\pi}}^{\alpha}\}$ are  of the first, second, $\ldots,$ up to the $\pi$th-orders. Denoting  \begin{equation}
D_i = \frac{ \partial }{ \partial x^i} + u_{i}^{\alpha } \frac{ \partial }{ \partial u^{\alpha }} + u_{ij}^{\alpha } \frac{ \partial }{ \partial u_{j}^{\alpha }}+ \ldots, \label{sec2.2.2}
\end{equation} the  total differentiation operator with respect to the variables $x^i$ and $ \delta_i^j$,  the Kronecker delta, we have
\begin{equation} \label{sec2.2.1} D_i(x^j)= \delta_i^j,\quad 
u_{i}^{\alpha}= D_i( u^{\alpha}), \quad u_{ij}^{\alpha}= D_{j}(D_{i}( u^{\alpha})),\ldots,
\end{equation} where variables $u^{\alpha}_{i} $ defined in (\ref{sec2.2.1}) are named as differential variables \citep{ibragimov}. A discussion of prolonged groups and prolonged generator follows.
\begin{enumerate}[(a)]
 \item \textbf{Prolonged groups} \newline
Consider the local Lie group $\mathcal{G}$ given by the transformations  \begin{align}
\begin{aligned}
\bar{x}^i = \varphi^i(x^i,u^{\alpha},\epsilon ), \quad \varphi^i\Big|_{\epsilon=0} = x^i,\quad 
\bar{u}^{\alpha } = \psi^{\alpha}(x^i,u^{\alpha},\epsilon),\quad \psi^{\alpha}\Big|_{\epsilon=0} = u^{\alpha },
\end{aligned} \label{sec2.3.1}
\end{align} where  the symbol $\Big|_{\epsilon=0} $ means evaluated on $\epsilon=0 $.
\begin{defn}
The construction of  the group $\mathcal{G}$ given by (\ref{sec2.3.1}) is an equivalence of  the computation of  infinitesimal transformations
\begin{equation} \bar{x}^i \approx x^i + \xi^i(x^i,u^{\alpha})  \epsilon, \quad \varphi^i\Big|_{\epsilon=0} = x^i,\quad \bar{u}^{\alpha } \approx u^{\alpha } + \eta^{\alpha} (x^i,u^{\alpha}) \epsilon,\quad \psi^{\alpha}\Big|_{\epsilon=0} = u^{\alpha }, \label{sec2.3.2}
\end{equation} obtained from (\ref{tranfmtn}) by a Taylor series expansion of  $ \varphi^i(x^i,u^{\alpha},\epsilon )$ and $ \psi^i(x^i,u^{\alpha},\epsilon )$ in $ \epsilon $ about $ \epsilon =0$ and keeping only the terms linear in $ \epsilon$, 
where \begin{equation} 
\xi^i(x^i,u^{\alpha}) = \frac{ \partial \varphi^i( x^i,u^{\alpha},\epsilon)}{ \partial \epsilon }\Big|_{\epsilon =0}, \quad  \eta^{\alpha}(x^i,u^{\alpha}) = \frac{ \partial \psi^{\alpha}( x^i,u^{\alpha},\epsilon)}{ \partial \epsilon }\Big|_{\epsilon =0}. \label{sec2.3.4}
\end{equation}
\end{defn}
\begin{rem}The symbol of infinitesimal transformations, $X$, is used to write (\ref{sec2.3.2}) as \begin{equation}
	\bar{x}^i \approx (1+X)x^i, \quad \bar{u}^{\alpha } \approx (1+X)u^{\alpha},\label{sec2.3.5}
	\end{equation} where  \begin{equation}
	X =  \xi^i(x^i,u^{\alpha}) \frac{ \partial }{ \partial x^i} + \eta^{\alpha}(x^i,u^{\alpha})  \frac{ \partial }{ \partial u^{\alpha }},	\label{sec2.3.6}
	\end{equation} is the generator  of the group $\mathcal{G}$ given by (\ref{sec2.3.1}).
\end{rem}

\begin{rem}
To obtain transformed derivatives from (\ref{tranfmtn}), we  use a change of variable formulae \begin{equation} D_i = D_i( \varphi^j) \bar{D}_j,
\label{sec2.3.7}
\end{equation} where $\bar{D}_j $ is the total differentiation in the  variables $ \bar{x}^i$. This means that \begin{equation} 
\bar{u}^{\alpha}_i =\bar{D}_i( \bar{u}^{\alpha}) , \quad \bar{u}^{\alpha}_{ij} = \bar{D}_j( \bar{u}^{\alpha}_i)= \bar{D}_i( \bar{u}^{\alpha}_j). \label{sec2.3.8}
\end{equation}
If we apply the change of variable formula given in (\ref{sec2.3.7}) on  $\mathcal{G}$ given by (\ref{sec2.3.1}), we get
\begin{equation} D_i( \psi^{\alpha}) = D_i( \varphi^j)\bar{D}_j ( \bar{u}^{\alpha} ) = \bar{u}^{\alpha}_j D_i( \varphi^j).
\label{sec2.3.9}
\end{equation}  Expansion of (\ref{sec2.3.9}) yields \begin{equation} \left( 
\frac{\partial \varphi^j}{\partial x^i}+ u_i^{\beta} \frac{\partial \varphi^j}{\partial u^{\beta}} \right) \bar{u}_j^{\beta} = \frac{\partial \psi^{\alpha}}{\partial x^i} + u_i^{\beta}  \frac{\partial \psi^{\alpha}}{\partial u^{\beta}}. \label{sec2.10.1}
\end{equation}	
The variables $ \bar{u}_i^{\alpha}$  can be written  as  functions of  $ x^i,u^{\alpha}, u_{(1)}$, that is
\begin{equation} \bar{u}_i^{\alpha} = \Phi^{\alpha}(x^i,u^{\alpha}, u_{(1)}, \epsilon), \quad \Phi^{\alpha}\Big|_{\epsilon =0} = u^{\alpha}_i. \label{sec2.3.11}
\end{equation}
\end{rem}
\begin{defn} \label{defproG}
The transformations in  the space of the  variables $ x^i,u^{\alpha}, u_{(1)} $ given in (\ref{sec2.3.1}) and (\ref{sec2.3.11}) form the first prolongation group $\mathcal{G}^{[1]}$.
\end{defn}
\begin{defn} Infinitesimal transformation of the first derivatives  is  \begin{equation}
\bar{u}_i^{\alpha} \approx  u_i^{\alpha} +  \zeta_i^{\alpha}\epsilon, \label{sec2.3.122}
\end{equation} where $ \zeta_i^{\alpha}=\zeta_i^{\alpha}(x^i,u^{\alpha},u_{(1)},\epsilon).$
\end{defn}
\begin{rem} In terms of infinitesimal transformations, the  first prolongation group $ \mathcal{G}^{[1]}$ is given  by (\ref{sec2.3.2}) and (\ref{sec2.3.122}).
\end{rem}
 \item \textbf{Prolonged generators} \newline
  A  description of how to extend a group generator is given.
\begin{defn}
By using the relation given in (\ref{sec2.3.9}) on the first prolongation group $ \mathcal{G}^{[1]}$  given by Definition \ref{defproG}, we obtain \citep{ibragimov2009practical}

\begin{align}
\label{p3} D_i( x^j +  \xi^j \epsilon)( u_j^{\alpha} +  \zeta_j^{\alpha}\epsilon )   =& D_i( u^{\alpha } +  \eta^{\alpha} \epsilon ), 
\end{align} 
which gives
\begin{align}
\label{p5} u^{\alpha}_i + \zeta_j^{\alpha} \epsilon+  u_j^{\alpha} \epsilon D_i \xi^j   =&  u^{\alpha }_i + D_i \eta^{\alpha} \epsilon ,
\end{align} 
and thus 
\begin{align}
\label{p44} \zeta_i^{\alpha} =& D_i(\eta^{\alpha}  )- u_j^{\alpha}  D_i (\xi^j),
\end{align}
is  the first prolongation formula.
\end{defn}
\begin{rem}
Similarly, we get  higher order prolongations \citep{ibragimov}, 
\begin{align}
\begin{aligned}
\zeta_{ij}^{\alpha} &= D_j(\zeta_i^{\alpha}  )- u_{i\kappa}^{\alpha}  D_j (\xi^{\kappa}),\\
\vdots  \\
\zeta_{i_1,\ldots,i_\kappa}^{\alpha} & = D_{i_\kappa}( \zeta_{i_1,\ldots,i_{\kappa-1}}^{\alpha} )- u_{i_1,i_2,\ldots, i_{\kappa-1}j}^{\alpha}  D_{i_\kappa} (\xi^{j}).
\end{aligned} \label{sec2.4.2}
\end{align}
\end{rem}
 \begin{rem}The  prolonged generators of the prolongations $ \mathcal{G}^{[1]},\ldots,  \mathcal{G}^{[\kappa]}$ of the group $\mathcal{G}$  are  
 	\begin{align} \begin{aligned}
 	X^{[1]} =& X + \zeta_i^{\alpha} \frac{ \partial }{ \partial u_i^{\alpha}},\\
 	\vdots \\
 	X^{[\kappa]} = &X^{[\kappa-1]} + \zeta_{i_1,\ldots,i_\kappa}^{\alpha} \frac{ \partial }{ \partial \zeta_{i_1,\ldots,i_\kappa}^{\alpha}}, \,\, \kappa \geq 1,
 	\end{aligned} \label{sec2.4.3}
 	\end{align} where $X$ is the group generator given by (\ref{sec2.3.6}).\end{rem}
\end{enumerate}
\subsection{ Group invariants}
This subsection  presents some basic definitions and a  theorem that are very important in Lie group analysis.
\begin{defn} A function $ \Gamma(x^i,u^{\alpha}) $ is called an invariant of  the group $\mathcal{G}$ of transformations given by (\ref{tranfmtn}) if 
	\begin{equation}
	\Gamma( \bar{x}^i,\bar{u}^{\alpha}) = \Gamma(x^i,u^{\alpha}). \label{sec2.6.1}
	\end{equation}  \end{defn}
\begin{thm} A function $\Gamma(x^i,u^{\alpha})$ is  an invariant of the group $\mathcal{G}$  given by (\ref{tranfmtn}) if and only if it solves the following first-order linear PDE: \citep{ibragimov} \begin{equation} X \Gamma = \xi^i(x^i,u^{\alpha}) \frac{ \partial \Gamma  }{ \partial x^i} + \eta^{\alpha}(x^i,u^{\alpha})  \frac{ \partial \Gamma }{ \partial u^{\alpha }} =0.\label{eqthr}
	\end{equation} \label{sec2.6.2}
\end{thm}
From Theorem (\ref{sec2.6.2}), we have the following result.
\begin{thm}
 The local Lie  group $\mathcal{G}$ of transformations in  $\mathbb{R}^n$ given by  (\ref{tranfmtn}) \citep{ibragimov2009practical} has  precisely $ n-1$ functionally independent invariants. One can take, as the basic invariants, the left-hand sides of the first integrals  \begin{equation}
	\psi_1(x^i,u^{\alpha})=c_1, \ldots,\psi_{n-1}(x^i,u^{\alpha}) = c_{n-1}, \label{sec2.6.3}
	\end{equation} of the characteristic equations for (\ref{eqthr}):
	\begin{equation}\frac{ \mathrm{d}x^i}{ \xi^i(x^i,u^{\alpha})}= \frac{ \mathrm{d}u^{\alpha}}{ \eta^{\alpha}(x^i,u^{\alpha})}. \label{sec2.6.4}
	\end{equation} \label{thm17}	
\end{thm}
\subsection{Symmetry groups }
This subsection presents some vital definitions on Lie groups of PDEs.  
\begin{defn} The vector field $X$ (\ref{sec2.3.6}) is  a Lie point symmetry of the  PDE system given in (\ref{pi-thorderpde}) if the  determining equations  
	\begin{equation}
	X^{[\pi]} \Delta_{ \alpha } \Big|_{ \Delta_{ \alpha }=0 } = 0, \quad \alpha=1,\ldots,m, \quad \pi \geq 1, \label{sec2.5.2}
	\end{equation} are satisfied, where  $\Big |_{\Delta_{ \alpha }=0}$ means evaluated on  $ \Delta_{\alpha} =0$ and $ 	X^{[\pi]}$ is the $\pi$-th prolongation of $X$.
\end{defn}
\begin{defn}The Lie group $\mathcal{G}$  is a symmetry group of the PDE system given in  (\ref{pi-thorderpde}) if the PDE system  (\ref{pi-thorderpde}) is  form-invariant, that is
	\begin{equation} \Delta_{\alpha} \left( \bar{x}^i, \bar{u}^{\alpha} , \bar{u}_{(1)}, \ldots, \bar{u}_{(\pi)}\right) =0. 
	\label{sec2.5.3}
	\end{equation}	
\end{defn}
\begin{thm} Given the infinitesimal  transformations in (\ref{sec2.3.1}), the Lie group  $\mathcal{G}$ in (\ref{tranfmtn}) is found by integrating the Lie equations
	\begin{align}
	\frac{ \mathrm{d} \bar{x}^i}{ \mathrm{d} \epsilon } = \xi^i(\bar{x}^i, \bar{u}^{\alpha}), \quad  \bar{x}^i\Big|_{\epsilon=0} = x^i,\quad
	\frac{ \mathrm{d}  \bar{u}^{\alpha} }{ \mathrm{d} \epsilon } = \eta^{\alpha}(\bar{x}^i, \bar{u}^{\alpha} ), \quad \bar{u}^{\alpha} \Big|_{ \epsilon=0} = u^{\alpha}.
	\end{align} \label{lieeqns}	
\end{thm}
\subsection{Lie algebras}
We use this subsection to define a Lie algebra of operators.
\begin{defn}   A vector space $ \mathcal{V}_{r}$ of operators \citep{ibragimov}  $X$ (\ref{sec2.3.6}) is  a Lie algebra if for any two operators, $X_i, X_j \in  \mathcal{V}_r$ , their commutator  \begin{equation} [X_i,X_j ]= X_iX_j-X_jX_i,\label{sec2.7.2} \end{equation} is in $ \mathcal{V}_r$  for all  $ i,j=1,\ldots,r$.
\end{defn}
\begin{rem}
The commutator satisfies the properties of bilinearity,  skew symmetry and  the Jacobi identity \citep{ibragimov}.	
\end{rem}
=======
\chapter{Groundwork for symmetr analysis and conservation laws}
This chapter is an exposition of fundamental concepts on  the theory of continuous groups, symmetry properties of partial  differential equations, conservation laws and exact solutions;  most of which  are adapted from \cite{ibragimov2009practical},\cite{bluman2013symmetries},\cite{bluman2008symmetry},\cite{olver2000applications},\cite{ovsyannikov2013lectures} \citep{arrigo2015symmetry}.
%% add books by IbrAGIMOV 
\section{One parameter groups of continuous transformations}
We consider $n$ independent variables $ x= ( x^1, \ldots, x^n ) $  with co-ordinates $ x^i$ and   $m$ dependent variables $ u= ( u^1, \ldots, u^m ) $ with co-ordinates $ u^{\alpha}$. Let $ T_{\epsilon}$ be the transformation:
\begin{equation}
T_{\epsilon}:\bar{x}^i = \varphi^i(x,u,\epsilon ), \quad \bar{u}^{\alpha } =\psi^{\alpha}(x,u,\epsilon) \label{2.1.1}
\end{equation} involving  a continuous real valued parameter $\epsilon$ in a neighbourhood $ D' \subset D \subset \mathbb{R}$ of $ \epsilon=0 $ for   $ \varphi^i$ and $ \psi^{\alpha}$ smooth functions of the variables $x^i$ and $u^{\alpha}$ and analytic in $\epsilon$.
\begin{defn} The set $\mathcal{G}$  of transformations (\ref{2.1.1}) is called a continuous one parameter local lie group of transformations in the space of variables $x$ and $u$ if:
	\begin{enumerate}[(i)]
		\item For $T_{\epsilon_1},T_{\epsilon_2}  \in \mathcal{G}$, where $\epsilon_1,\epsilon_2 \in  D' \subset D $, then $T_{\epsilon_1}T_{\epsilon_2} = T_{\epsilon_3} \in \mathcal{G}, \epsilon_3= \phi( \epsilon_1,\epsilon_2)$	(Closure)	
		\item $  T_0 \in \mathcal{G} $ if and only if $ \epsilon=0$ such that $ T_{\epsilon}T_{0}=T_{0}T_{\epsilon}=T_{\epsilon}$(Identity)
		\item For $ T_{\epsilon} \in \mathcal{G}, \epsilon\in  D' \subset D$, there exists a unique $ T_{\epsilon^{-1}} \in \mathcal{G}, \epsilon^{-1} \in D $
		such that $ T_{\epsilon} T_{\epsilon^{-1}} =T_{\epsilon^{-1}}T_{\epsilon}=T_{0}$(Inverse)
	\end{enumerate}

\begin{rem}
 Associativity property of the group $\mathcal{G} $ follows from (i). \label{2.1}
\end{rem}
\end{defn}
\section{Prolongations}
One of our major interests is to  describe how the derivatives of $u^{\alpha}$ transform. The  operator for total differentiation, $D_i$ is defined as
\begin{equation}
D_i = \frac{ \partial }{ \partial x^i} + u_{i}^{\alpha } \frac{ \partial }{ \partial u^{\alpha }} + u_{ij}^{\alpha } \frac{ \partial }{ \partial u_{j}^{\alpha }}+ \ldots \quad i,j = 1, \ldots, n, \quad \alpha=1, \ldots, m. \label{2.2.2}
\end{equation}
The following derivatives of $x$ and $u$ are written as \begin{equation} \label{2.2.1} D_i(x^j)= \delta_i^j,\quad 
u_{i}^{\alpha}= D_i( u^{\alpha}), \quad u_{ij}^{\alpha}= D_{j}(D_{i}( u^{\alpha})), \quad \ldots 
\end{equation} where $ \delta_i^j$ is the Kronecker delta symbol, that is,
\begin{align}
	\delta_i^j = \begin{cases} 1,\quad  \text{if} \quad i=j\\ 0 \quad \text{if} \quad i \neq j
	\end{cases}
\end{align}
The collection of all first derivatives $ u_i^{\alpha}$ is denoted by $ u_{(1)} $, namely, \begin{equation}
u_{(1)} = \{ u_i^{\alpha}  \}, \quad \alpha = 1, \ldots, m, \quad i =1, \ldots, n
\end{equation}
The second and higher derivatives  up to $ \pi$th-order derivative are denoted analogously, namely, \begin{equation}
 u_{(2)}=  \{ u_{ij}^{\alpha}  \}, \ldots, u_{(\pi)} = \{ u_{i_1, \ldots i_\pi}^{\alpha}  \}
\end{equation}


\begin{rem}
We want all the variables $ x, u , u_{(1)}, \ldots $ to be functionally independent and connected only by the differential  relations (\ref{2.2.1})  \citep{ibragimov1999elementary}.

The  variables $u^{\alpha}_{i} $ are called \textbf{differential variables.}

The variables  $ u_{ij}, \ldots $ are assumed to be symmetric in subscripts, that is $u_{ij} = u_{ji}, \ldots$
\end{rem}
\subsection{ Prolonged groups}
Let $ z= (x,u)$ and one parameter group of transformations $\mathcal{G}$  be \begin{align}
\begin{aligned}
\bar{x}^i = \varphi^i(x,u,\epsilon ), \quad \varphi^i|_{\epsilon=0} = x^i, \quad i =1,\ldots,n\\
\bar{u}^{\alpha } =\psi^{\alpha}(x,u,\epsilon),\quad \psi^{\alpha}|_{\epsilon=0} = u^{\alpha }, \quad \alpha =1,\ldots,m
\end{aligned} \label{2.3.1}
\end{align} where the symbol $|_{\epsilon=0} $ means evaluated on $\epsilon=0 $.

According to Lie's theory, the construction of the group $\mathcal{G}$ is the same as determination of corresponding infinitesimal transformations:
\begin{equation} \bar{x}^i \approx x^i + \epsilon \xi^i(x,u), \quad \bar{u}^{\alpha } \approx u^{\alpha } + \epsilon  \eta^{\alpha} ( x,u) \label{2.3.2}
\end{equation} obtained from (\ref{2.1.1}) by Taylor series expansion of the functions $ \varphi^i(x,u,\epsilon )$ and $ \psi^i(x,u,\epsilon )$ in $ \epsilon $ about $ \epsilon =0$ while taking into account the initial conditions \begin{equation}
\varphi^i|_{\epsilon=0} = x^i, \quad \psi^{\alpha}|_{\epsilon=0} = u^{\alpha } \label{2.3.3}
\end{equation} and keeping only the terms linear in $ \epsilon$. As a result, we have \begin{equation} 
\xi^i(x,u) = \frac{ \partial \varphi^i( x,u,\epsilon)}{ \partial \epsilon }|_{\epsilon =0}, \quad  \eta^{\alpha}(x,u) = \frac{ \partial \psi^{\alpha}( x,u,\epsilon)}{ \partial \epsilon }|_{\epsilon =0} \label{2.3.4}
\end{equation}
We use a \textit{symbol} $ X$  of infinitesimal transformations  to write (\ref{2.3.2}) as \begin{equation}
\bar{x}^i \approx (1+X)x, \quad \bar{u}^{\alpha } \approx (1+X)u\label{2.3.5}
\end{equation} where the differential operator \begin{equation}
X =  \xi^i(x,u) \frac{ \partial }{ \partial x^i} + \eta^{\alpha}(x,u)  \frac{ \partial }{ \partial u^{\alpha }} 	\label{2.3.6}
\end{equation} is known as the \textbf{infinitesimal operator or generator }of the group $\mathcal{G}$. We say that $X$ is an \textit{admitted operator} or an \textit{infinitesimal symmetry} of equation(\ref{2.2.4}) if the group $\mathcal{G}$ is admitted by equation (\ref{2.2.4}).
 Transformed derivatives are obtained from (\ref{2.1.1}) by use of  change of variable formulae \begin{equation} D_i = D_i( \varphi^j) \bar{D}_j
\label{2.3.7}
\end{equation} where $\bar{D}_j $ is the total differentiation in transformed variables $ \bar{x}^i$. This means that \begin{equation} 
\bar{u}^{\alpha}_i =\bar{D}_i( \bar{u}^{\alpha}) , \quad \bar{u}^{\alpha}_{ij} = \bar{D}_j( \bar{u}^{\alpha}_i)= \bar{D}_i( \bar{u}^{\alpha}_j) \label{2.3.8}
\end{equation}
By applying (\ref{2.3.7}) on the group $G$ (\ref{2.3.1}), we have 
\begin{equation} D_i( \psi^{\alpha}) = D_i( \varphi^j)\bar{D}_j ( \bar{u}^{\alpha} ) = \bar{u}^{\alpha}_j D_i( \varphi^j)
\label{2.3.9}
\end{equation}  or \begin{equation} \left( 
\frac{\partial \varphi^j}{\partial x^i}+ u_i^{\beta} \frac{\partial \varphi^j}{\partial u^{\beta}} \right) \bar{u}_j^{\beta} = \frac{\partial \psi^{\alpha}}{\partial x^i} + u_i^{\beta}  \frac{\partial \psi^{\alpha}}{\partial u^{\beta}} \quad i,j = 1,\ldots,n, \quad \alpha,\beta =1,\ldots,m. \label{2.10.1}
\end{equation}

One  can represent the quantities $ \bar{u}_i^{\alpha}$ as functions of $ x, u, u_{(1)}$, that is, 
\begin{equation} \bar{u}_i^{\alpha} = \Phi^{\alpha}(x,u, u_{(1)}, \epsilon), \quad \Phi^{\alpha}|_{\epsilon =0} = u^{\alpha}_i \label{2.3.11}
\end{equation}
The transformations in the space of variables $ x,u, u_{(1)} $ given by (\ref{2.3.1}) and (\ref{2.3.11}) form a one-parameter group called the \textit{first prolongation} or \textit{extension} of the group $\mathcal{G}$ denoted by $\mathcal{G}^{[1]}$.

The infinitesimal transformation of the first derivatives $ u_i^{\alpha}$ is given as \begin{equation}
\bar{u}_i^{\alpha} \approx  u_i^{\alpha} + \epsilon \zeta_i^{\alpha}  \label{2.3.12}
\end{equation} so that the first prolongation group $ \mathcal{G}^{[1]}$ is (\ref{2.3.2}) and (\ref{2.3.12}). Higher-order prolongations of $\mathcal{G}$ , namely, $\mathcal{G}^{[2]}, \mathcal{G}^{[3]},
\ldots$ can be obtained  analogously by taking successive derivatives of (\ref{2.3.9}).
\subsection{Prolonged generators}
By using (\ref{2.3.9}) on the  group $ \mathcal{G}^{[1]}$, we obtain 
\begin{align}
\begin{aligned}
D_i( \varphi^j)( \bar{u}^{\alpha}_j )& =   D_i( \psi^{\alpha})\\
D_i( x^j + \epsilon \xi^j )( u_j^{\alpha} + \epsilon \zeta_j^{\alpha} )& =   D_i( u^{\alpha } + \epsilon  \eta^{\alpha} ) \\
(\delta^i_j + \epsilon D_i\xi^j )( u_j^{\alpha} + \epsilon \zeta_j^{\alpha} ) &=  u^{\alpha }_i + \epsilon D_i \eta^{\alpha} \\
u^{\alpha}_i +\epsilon \zeta_j^{\alpha} + \epsilon u_j^{\alpha}  D_i \xi^j  &=  u^{\alpha }_i + \epsilon D_i \eta^{\alpha} \\
\zeta_i^{\alpha}&= D_i(\eta^{\alpha}  )- u_j^{\alpha}  D_i (\xi^j) \quad \text{(sum on $j$)},\quad i,j =1,\ldots,n
\end{aligned} \label{2.4.1}
\end{align}
The last term of (\ref{2.4.1}) is called the \textit{first prolongation formula}. Analogously \citep{ibragimov1999elementary}, one can recursively have 
\begin{align}
\begin{aligned}
\zeta_{ij}^{\alpha} &= D_j(\zeta_i^{\alpha}  )- u_{i\kappa}^{\alpha}  D_j (\xi^{\kappa}) \quad   (\text{sum on $ \kappa$})\\
\vdots  \\
\zeta_{i_1,\ldots,i_\kappa}^{\alpha} & = D_{i_\kappa}( \zeta_{i_1,\ldots,i_{\kappa-1}}^{\alpha} )- u_{i_1,i_2,\ldots, i_{\kappa-1}j}^{\alpha}  D_{i_\kappa} (\xi^{j}) \quad   (\text{sum on  $j$})
\end{aligned} \label{2.4.2}
\end{align}
The prolongations of the group $\mathcal{G}$ form a group denoted by $ \mathcal{G}^{[1]},\ldots,  \mathcal{G}^{[\kappa]}$. 
The corresponding prolonged generators are 
\begin{align} \begin{aligned}
X^{[1]} =& X + \zeta_i^{\alpha} \frac{ \partial }{ \partial u^i_{\alpha}} \quad (\text{sum on $ i, \alpha $ } ) \\
\vdots \\
X^{[\kappa]} = &X^{[\kappa-1]} + \zeta_{i_1,\ldots,i_\kappa}^{\alpha} \frac{ \partial }{ \partial \zeta_{i_1,\ldots,i_\kappa}^{\alpha}} \quad \kappa \geq 1
\end{aligned} \label{2.4.3}
\end{align}, where $X$ is as given in (\ref{2.3.6}).

\section{ Group invariants}
We define invariance and mention  a key theorem on necessary and sufficient condition for invariance of a function.
\begin{defn} A function $ \Gamma(x,u) $ is said to be an \textit{invariant }of the group  $\mathcal{G}$ (\ref{2.1.1}), if 
	\begin{equation}
	\Gamma( \bar{x},\bar{u} ) \equiv \Gamma ( \varphi^i(x,u,\epsilon), \psi^{\alpha}(x,u,\epsilon))= \Gamma(x,u) \label{2.6.1}
	\end{equation} identically in $ x, u$ and $ \epsilon$.
\end{defn}


\begin{thm} A necessary and sufficient condition for a  function $\Gamma(x,u)$ to be an invariant under the symmetry $X$(\ref{2.3.5}) is that, \begin{equation} X \Gamma = \xi^i(x,u) \frac{ \partial \Gamma  }{ \partial x^i} + \eta^{\alpha}(x,u)  \frac{ \partial \Gamma }{ \partial u^{\alpha }} =0,\quad i = 1,  \ldots, n, \quad \alpha = 1,  \ldots, m 
	\end{equation} \label{2.6.2}
\end{thm}
\begin{rem}
	It follows from theorem (\ref{2.6.2}), that every one parameter group of point transformations (\ref{2.1.1}) has $ n-1$ functionally independent invariants, which can be taken as the left-hand side of any first integrals \begin{equation}
	J_1(x,u)=C_1, \ldots, J_{n-1}(x,u) = C_{n-1} \label{2.6.3}
	\end{equation} of the characteristic equations 
	\begin{equation}\frac{ \mathrm{d}x^1}{ \xi^1(x,u)}= \ldots= \frac{ \mathrm{d}x^n}{ \xi^n(x,u)}= \frac{ \mathrm{d}u^1}{ \eta^1(x,u)}=\ldots = \frac{ \mathrm{d}u^m}{ \eta^m(x,u)}. \label{2.6.4}
	\end{equation}	
\end{rem}



\section{Symmetry groups of  partial differential equations}
Here we present some vital definitions on Lie groups admitted by partial differential equations. 

\begin{rem}The definitions of one-parameter groups and symmetry groups of partial differential equations can be extended from those of differential equations \citep*{ibragimov2009practical}.
\end{rem} 
We shall consider a system of  $ \pi$th-order partial differential equations, namely,
\begin{equation}
E_{\alpha}( x, u , u_{(1)}, \ldots,u_{(\pi)}) =0, \quad \alpha = 1,  \ldots, m. \label{2.2.4}
\end{equation}

\begin{defn} We say that the vector field $X$ (\ref{2.3.6}) is  a \textbf{point symmetry} of equation ( \ref{2.2.4}) if the determining equation (\ref{2.5.2}) is satisfied
	\begin{equation}
	X^{[\pi]} \left(  E_{ \alpha } \textbar_{ E_{ \alpha }=0 } \right)= 0, \quad \alpha=1,\ldots,m \label{2.5.2}
	\end{equation}	where the symbol $\textbar_{E_{ \alpha }=0}$ means evaluated on the equation $ E_{\alpha} =0$.
\end{defn}
\begin{rem}We refer to  equation (\ref{2.5.2}) as  the \textbf{determining equation} of equation (\ref{2.2.4}) because it determines all infinitesimal symmetries of equation (\ref{2.2.4}).
\end{rem}

\begin{defn}We shall to refer to a  one-parameter group $\mathcal{G}$ of transformations as a \textbf{symmetry group} of equation (\ref{2.2.4}) if equation (\ref{2.2.4}) is form-invariant in the new variables $ \bar{x}$ and $ \bar{u}$, namely,
	\begin{equation} E_{\alpha}( \bar{x}, \bar{u} , \bar{u}_{(1)}, \ldots, \bar{u}_{(\pi)})=0 
	\label{2.5.3}
	\end{equation} for the function $ E_{\alpha}$  same as in equation (\ref{2.2.4}).	
\end{defn}
\begin{thm} If we have the infinitesimal  transformation (\ref{2.3.1})  or its symbol $X$, then the corresponding one-parameter local lie group  $\mathcal{G}$ is obtained by integrating the \textbf{Lie equations }
	\begin{align}
	\begin{aligned}
	\frac{ \mathrm{d} \bar{x}^i}{ \mathrm{d} \epsilon } = \xi^i(\bar{x}, \bar{u} ), \quad  \bar{x}^i|_{\epsilon=0} = x,\\
	\frac{ \mathrm{d}  \bar{u}^{\alpha} }{ \mathrm{d} \epsilon } = \eta^{\alpha}(\bar{x}, \bar{u} ), \quad \bar{u}^{\alpha} |_{ \epsilon=0} = u
	\end{aligned}
	\end{align} 	
\end{thm}





\section{Lie algebras}
This section is dedicated to definitions involving Lie algebra of operators. 
\begin{defn} A vector space $ \mathcal{V}_{r}$ of operators \citep{ibragimov1999elementary}  $X$ (\ref{2.3.6}) is said to be a \textbf{Lie algebra} if for any two operators, $X_i, X_j \in  \mathcal{V}_r$ , then the \textbf{commutator} $[X_i,X_j ]$ defined as, \begin{equation} [X_i,X_j ]= X_iX_j-X_jX_i, \label{2.7.2}\end{equation} is also in $ \mathcal{V}_r$  for $ i,j=1,\ldots,r$.
\end{defn}
\begin{rem}
It follows  from definition (\ref{2.7.2}) that the commutator is bilinear, skew-symmetric and satisfies the Jacobi identity.	
\end{rem}

>>>>>>> f3e458080f41e5290173f14b8121103bb6a456e0
\begin{thm}
The set of  solutions of the  determining equation given by  (\ref{sec2.5.2}) forms a Lie algebra\citep{ibragimov}. 
\end{thm}
<<<<<<< HEAD
\section{Conservation laws}
This section serves as an introduction to conservation laws and the algorithms used to find them, that is, the method of multipliers and Ibragimov's conservation theorem. We will use these methods in the next part of the dissertation.
\subsection{Fundamental operators}
Let a system of $\pi$th-order PDEs be given by (\ref{pi-thorderpde}).
\begin{defn} The Euler-Lagrange operator ${\delta}/{\delta u^{\alpha}}$ is 
=======
\begin{rem}
One can always define subalgebra and ideal for lie algebras as in done for groups.
\end{rem}
\section{Exact solutions of non-linear partial differential equations}
A quick flavour of some three methods for obtaining  exact solutions of partial differential equations is given.
\subsection{(G'/G)-expansion method }
%% wang etal[8]
\citep{wang2008g}
The function $G$ satisfies a second order linear ordinary differential equation with constant coefficients. The method presents travelling wave solutions expressible as trigonometric, rational or hyperbolic functions. Consider a non linear partial differential equation of the form
 \begin{equation}
E_{1}(u,u_{t}, u_{x}, u_{tt},u_{xx} \ldots) =0,  \label{2.8.1}
\end{equation} The equation(\ref{2.8.1}) is first transformed into a non linear ODE through the substitution
\begin{equation}
u(t,x) = \nu(\varPhi), \quad \varPhi = x- \rho t, \label{2.8.2}
\end{equation} where $\rho$ is a constant, which yields
\begin{equation} E_{1}(\nu,-\rho \nu', \nu',\rho^2\nu'', -\rho \nu'',\nu'' \ldots) =0,  \label{2.8.3}
\end{equation}
Then assuming that travelling wave solutions of \ref{2.8.3} can be expressed as a polynomial in $\frac{ G'(\varPhi)}{G(\varPhi)} $ \begin{equation} \nu(\varPhi) = \sum_{i=0}^{m} \mathcal{A}_i \left( \frac{ G'(\varPhi)}{G(\varPhi)}\right)^i  \label{2.8.4} 
\end{equation} where $ \mathcal{A}_0, \ldots, \mathcal{A}_m$ are to be determined. Determination of  $m$ involves considering the homogeneous balance between the highest derivatives and non-linear terms appearing equation (\ref{2.8.3}). In  equation \ref{2.8.4}, $ G(\varPhi)$ satisfies the second order  linear homogeneous ordinary differential equation \begin{equation} G''(\varPhi) + \lambda G'(\varPhi) + \mu G(\varPhi) =0 \label{2.8.5}
\end{equation} for some arbitrary contants $ \lambda$ and $\mu$.
\begin{rem}
The substitution of $\nu(\varPhi)$ from equations (\ref{2.8.4}) and (\ref{2.8.5}) into equation (\ref{2.8.3}) yields a polynomial in $(G'/G)^i$, whose coefficients are equated to zero and the resultant system of algebraic equations solved for the values of $A_i$.
\end{rem}
\subsection{Extended Jacobi elliptic function expansion method}
% Give a brief history of the method, who brought it
%[28-31]

The steps involved in this technique are explained below \citep{hong2013new},\citep{hong2009extended}.
The equation(\ref{2.8.1}) is first transformed into the ordinary differential equation (\ref{2.8.3}) just like in the $(G'/G)$ expansion method after which the  solutions $\nu(\varPhi) $ are assumed to be expressible in the form
\begin{equation}
 \nu(\varPhi) = \sum_{i=-\mathcal{M}}^{\mathcal{M}}A_i \mathcal{H}(\varPhi)^i  \label{2.8.6} 
\end{equation} for some positive integer, $\mathcal{M}$ obtainable by the balancing procedure and \begin{equation} \mathcal{H}(\varPhi) = cn(\varPhi|\omega), \label{2.8.7}
\end{equation} the cosine amplitude function satisfying the first order ordinary differential equation %%[28,32]
\begin{equation} \mathcal{H}'(\varPhi) = - \sqrt{ \left(  1- \mathcal{H}^2(\varPhi) \right) \left( 1- \omega + \omega \mathcal{H}^2(\varPhi) \right)  }, \label{2.8.8}
\end{equation} and the sine amplitude function \begin{equation}
	\mathcal{H}(z) = sn( \varPhi|\omega) \label{2.8.9}
\end{equation} satisfying the first order ordinary differential equation 
\begin{equation}
	\mathcal{H}'(\varPhi) = - \sqrt{ \left(  1- \mathcal{H}^2(\varPhi) \right) \left( 1-\omega \mathcal{H}^2(\varPhi) \right)  }, \label{2.8.10}
\end{equation}
\begin{rem} The value of $ \nu(\phi)$ from equation (\ref{2.8.6}) is then substituted into either of equations (\ref{2.8.8}) and (\ref{2.8.10}) to obtain a polynomial in powers of $ \mathcal{H}(\phi)$, again whose coefficients are separated resulting into an algebraic system of equations which are then solved for the values of $ A_i, \quad i = 0,\pm 1,\ldots,\pm \mathcal{M}$.
\end{rem}
\subsection{Kudryashov method}
%%%  reference [9] ,[35,36] This method was introduced by 
We  consider a non-linear partial differential equation of the form
\begin{equation}
E_{2}(u,u_{t}, u_{x}, u_{tt},u_{xx} \ldots) =0,  \label{2.8.11}
\end{equation}
The iterations involved in the algorithm are as follows  \citep{kudryashov2012one}.
\begin{enumerate}[o]
\item A substitution $u(t,x) = \nu(\varPhi), \quad \varPhi =  \tau x+\omega t$ for $\tau$ and $\omega$  constants transforms the equation (\ref{2.8.11}) to the ordinary differential equation
\begin{equation} E_{1}(\nu,\omega \nu', \tau \nu',\omega^2\nu'', \tau^2 \nu'', \ldots) =0,  \label{2.8.13}
\end{equation}
\item Supposing that the exact solutions of equation (\ref{2.8.13}) is expressible as \begin{equation}
\nu(\varPhi) = \sum_{ n=0}^{N} a_n \textit{Q}^n(\varPhi) \label{2.8.14}
\end{equation} wherein $a_n$ are constants to be determined for $a_N \neq 0$, then $\textit{Q}(\varPhi)$ satisfies the first order non-linear ordinary differential equation \begin{equation}\textit{Q}_{\varPhi} ( \varPhi) =\textit{Q}^2(\varPhi) -\textit{Q}(\varPhi). \label{2.8.15}
\end{equation}
The solution of equation (\ref{2.8.15}) is given as \begin{equation}\textit{Q}(\varPhi) = \frac{1}{ 1 + e^{\varPhi}} \label{2.8.16}
\end{equation}
\item Then the value of $\nu(\varPhi)$ is substituted into equation (\ref{2.8.13})	 and equation (\ref{2.8.15}) used to generate an equation which involves the powers of $\textit{Q}$.
\item Now equating the powers of $\textit{Q}$ to zero outputs the system of algebraic equation 
\begin{equation}\mathcal{P}_n ( a_N,a_{N-1},\ldots,a_0,\tau,\omega , \ldots) =0, \quad n= 0, \ldots, N \label{2.8.17}
\end{equation}
\end{enumerate}
\begin{rem} Last but not least, the solution of algebraic equations in (\ref{2.8.17}) gives values of the coefficients $a_0,\ldots,a_N$ alongside relations for parameters of equation (\ref{2.8.13}). Ultimately, exact solutions of equation (\ref{2.8.13})  expressed in the form in equation (\ref{2.8.14}) are obtained.
\end{rem}
\section{Conservation laws}
This is the idea that a quantity of a physical system does not dissipate with  time evolution. We discuss  briskly fundamental operators, Noether's theorem and alternative techniques for  deriving  conservation laws for systems that do have admit a Lagrangian.
\subsection{Fundamental operators and their relationship}
Let a system of $\pi$th-order partial differential equations with $n$ independent variables and $ x= (x_1, \ldots, x^n)$  and  $m$ dependent variables $ u= ( u^1, \ldots, u^m ) $ be given by equation (  \ref{2.2.4}).
\begin{defn} The Euler-Lagrange operator , for each $\alpha$ is defined by,
>>>>>>> f3e458080f41e5290173f14b8121103bb6a456e0
	\begin{equation}
	\frac{ \delta}{ \delta u^{\alpha}} = \frac{ \partial }{ \partial  u^{\alpha}} + \sum_{\substack{  \kappa \geq 1}} (-1)^{\kappa} D_{   i_1}, \ldots, D_{   i_{\kappa}} \frac{ \partial }{ \partial u^{\alpha }_{ i_1i_2\ldots i_{\kappa }}},  \label{sec2.9.1}
	\end{equation} and the  Lie- B\"acklund operator in abbreviated form  \citep{ibragimov} is  \begin{equation}
		X =  \xi^i \frac{ \partial }{ \partial x^i} + \eta^{\alpha}   \frac{ \partial }{ \partial u^{\alpha }}+\ldots. \label{sec2.9.3}
		\end{equation}  
\end{defn}
<<<<<<< HEAD
\begin{rem}The Lie- B\"acklund operator  (\ref{sec2.9.3}) in its prolonged form is \begin{equation}
	X =  \xi^i \frac{ \partial }{ \partial x^i} + \eta^{\alpha}   \frac{ \partial }{ \partial u^{\alpha }}+\sum_{\substack{ \kappa  \geq 1}} \zeta_{i_1\ldots i_{\kappa}} \frac{ \partial }{ \partial u^{\alpha }_{ i_1i_2\ldots i_\kappa}}, \label{sec2.9.4}
	\end{equation} where  \begin{align}
	\zeta_{i}^{\alpha} &= D_i(W^{\alpha})+ \xi^{j} u_{i j }^{\alpha}, \quad j=1,\ldots,n,\\
	\vdots\\
	\zeta_{i_1\ldots i_{\kappa}}^{\alpha}&= D_{i_1 \ldots i_{\kappa} } (W^{\alpha}) + \xi^{j} u_{ j i_1 \ldots i_{\kappa} }^{\alpha}, \,\,  \end{align}
	and  the Lie characteristic function is \begin{align}	W^{\alpha} = \eta^{\alpha} -\xi^j u^{\alpha}_{j}. \label{sec2.9.5}
\end{align} 
\end{rem}
\begin{rem} The characteristic form of Lie- B\"acklund operator (\ref{sec2.9.4}) is \begin{equation}
	X =  \xi^i D_i  + W^{\alpha} \frac{ \partial }{ \partial u^{\alpha}}  +  D_{i_1 \ldots i_{\kappa} } (W^{\alpha})  \frac{ \partial }{ \partial u^{\alpha }_{ i_1i_2\ldots i_\kappa}}.  \label{sec2.9.6}
=======

\begin{defn} If there exits a function $ \mathcal{L} =\mathcal{L}( x, u , u_{(1)}, \ldots,u_{(s)} ),\quad s\leq \pi $ for $\pi$th-order of equation ( \ref{2.2.4}) such that the variational derivative of $\mathcal{L}$ vanishes, that is,
	\begin{equation}
	\frac{ \delta \mathcal{L}}{ \delta u^{\alpha}} =0, \alpha = 1, \ldots,m.  \label{2.9.2}
	\end{equation} then $ \mathcal{L}$ is called a Lagrangian of equation (\ref{2.2.4}). The equation ( \ref{2.9.2}) is called the Euler-Lagrange equation.
	
\end{defn}

\begin{defn}  The Lie- B\"acklund operator is given by \begin{equation}
	X =  \xi^i \frac{ \partial }{ \partial x^i} + \eta^{\alpha}   \frac{ \partial }{ \partial u^{\alpha }}, \quad  \xi^i ,\eta^{\alpha} \in \mathcal{A} \label{2.9.3}
	\end{equation} for some $ \mathcal{A} $ the space of differential functions \citep{ibragimov1999elementary}. 
\end{defn}
The operator  (\ref{2.9.3}) in its prolonged form is \begin{equation}
X =  \xi^i \frac{ \partial }{ \partial x^i} + \eta^{\alpha}   \frac{ \partial }{ \partial u^{\alpha }}+\sum_{\substack{ \kappa  \geq 1}} \zeta_{i_1\ldots i_{\kappa}} \frac{ \partial }{ \partial u^{\alpha }_{ i_1i_2\ldots i_\kappa}}, \label{2.9.4}
\end{equation} with the additional terms uniquely determined by the prolongation formulae \begin{align} \begin{aligned}
\zeta_{i}^{\alpha} &= D_i(W^{\alpha})+ \xi^{\iota} u_{i \iota }^{\alpha}, \quad \iota=1,\ldots,n\\
\vdots\\
\zeta_{i_1\ldots i_{\kappa}}^{\alpha}&= D_{i_1 \ldots i_{\kappa} } (W^{\alpha}) + \xi^{\iota} u_{ \iota i_1 \ldots i_{\kappa} }^{\alpha}, \kappa >1,\quad \iota=1,\ldots,n,\quad
 \alpha=1,\ldots,m.   
\end{aligned} \label{2.9.5}
\end{align}
for which  $ W^{\alpha} = \eta^{\alpha} -\xi^j u^{\alpha}_{\iota}$ (\ref{2.9.5}) is the \textit{Lie characteristic function}. 
The  Lie- B\"acklund operator (\ref{2.9.3})  can be written in characteristic form as \begin{equation}
X =  \xi^i D_i  + W^{\alpha} \frac{ \partial }{ \partial u^{\alpha}}  +  D_{i_1 \ldots i_{\kappa} } (W^{\alpha})  \frac{ \partial }{ \partial u^{\alpha }_{ i_1i_2\ldots i_\kappa}}   \label{2.9.6}
\end{equation} 
\section*{Noether's theorem on conservation laws}
\begin{defn} 
	The  \textit{Noether operators} associated with a Lie-B\"acklund symmetry  operator  $X$ are given by \begin{equation} N^i = \xi^i + W^{\alpha} \frac{ \delta }{ \delta u_i^{\alpha}} + \sum_{\substack{ \kappa  \geq 1}} D_{i_1 \ldots i_{\kappa} } (W^{\alpha})  \frac{ \delta }{ \delta u^{\alpha }_{ i_1i_2\ldots i_\kappa}}, \quad i= 1, \ldots, n  \label{2.9.7}
>>>>>>> f3e458080f41e5290173f14b8121103bb6a456e0
	\end{equation}
\end{rem}
\subsection{The method of multipliers}
\begin{defn}
 A function    $  \Lambda^{\alpha} \left( x^i, u^{\alpha},u_{(1)}, \ldots \right) =\Lambda^{\alpha},$  is a multiplier of the PDE system given by (\ref{pi-thorderpde}) if it satisfies the condition that \citep{olver2000applications}
\begin{equation}  \Lambda^{\alpha} \Delta_{\alpha} = D_iT^i, \label{sec2.9.13} 
\end{equation} where $D_iT^i$ is a divergence expression. \label{def2.2.6}
\end{defn}
\begin{defn}	To find the multipliers $\Lambda^{\alpha}$, one solves the determining equations (\ref{sec2.9.14})  \citep{bluman2008symmetry}, \begin{equation} \frac{ \delta }{ \delta u^{\alpha}} ( \Lambda^{\alpha} \Delta_{\alpha} )=0. \label{sec2.9.14}
\end{equation}
<<<<<<< HEAD
\end{defn}

\subsection{Ibragimov's conservation theorem \label{ibra}}
The technique \citep{ibragimov2007new}  enables one to construct conserved vectors associated with each  Lie point symmetry of the PDE system given by (\ref{pi-thorderpde}).
 \begin{defn} The adjoint equations of the system given by (\ref{pi-thorderpde})  are  \begin{equation} \Delta_{\alpha}^{*}\left( x^i, u^{\alpha} , v^{\alpha}, \ldots,u_{(\pi)},v_{(\pi)} \right) \equiv   \frac{ \delta }{ \delta u^{\alpha }} ( v^{\beta} \Delta_{\beta})   =0, \label{sec2.9.15}
	\end{equation} where  $ v^{\alpha}$ is the new dependent variable.
=======

The Lie-B\"acklund, Euler-Lagrange and Noether operators are  related by operator identity \citep{ibragimov1999elementary}\begin{equation} X + D_i( \xi^i) = W^{\alpha} \frac{ \delta }{ \delta u^{\alpha}} + D_i N^i,\quad  \quad i = 1, \ldots, n, \quad  \alpha = 1, \ldots,m. \label{2.9.9}
\end{equation}

\begin{defn} A lie B\"acklund operator $X$ of the form (\ref{2.9.3}) is called a\textit{ Noether symmetry} corresponding to  a Lagrangian $ \mathcal{L} \in \mathcal{A} $ if there exists a vector $ B^i = ( B^1, \ldots, B^n), B^i \in \mathcal{A}$ such that \begin{equation} X(\mathcal{L}) + \mathcal{L} D_i( \xi^i) = D_i(B^i) \label{2.9.10}
	\end{equation} if $ B^i=0,\quad ( i = 1,\ldots,n)$, then $X$ is called a \textit{strict Noether symmetry }corresponding to a Lagrangian $ \mathcal{L} \in \mathcal{A}$.
\end{defn}
\begin{thm}  Every Noether symmetry generator $X$ associated to a given Lagrangian $ \mathcal{L} \in \mathcal{A}$  of ( \ref{2.2.4}) , has a  corresponding conserved vector $ T = ( T^1, \ldots, T^n), T^i \in \mathcal{A} $, given by 
	\begin{equation}
	T^i = N^i( \mathcal{L})-B^i, \quad i=1, \ldots, n, \label{0}
	\end{equation} satisfying \begin{equation} D_i T^i|_{(\ref{2.2.4})}=0
	\label{2.9.11}	\end{equation} where the equation (\ref{2.9.11}) defines a\textbf{ local conservation law} for the system (\ref{2.2.4}).
\end{thm}
\begin{rem}
	In this Noether approach, the Lagrangian $\mathcal{L}$ is found and then equation (\ref{2.9.10}) is used to unleash Noether symmetries after which the equation (\ref{0}) will birth the corressponding Noether conserved vectors.
\end{rem}
\subsection{The multiplier method}
% a book should be cited here.
\cite{naz2012conservation}
This technique reduces the calculation of conservation laws to just solving a system of linear determining equations analogous to that of calculating symmetries. It explores the possibility of bringing an equation into conservation law by multiplying by  a specific function,  $ \Delta^{\alpha} ( x, u,u_{(1)}, \ldots)$,  called a \textit{multiplier}, for which the following must be satisfied in case of the partial differential equation system (\ref{2.2.4})  (\cite{olver2000applications},\citep{bluman2010applications})
\begin{equation}  \Delta^{\alpha} E_{\alpha} = D_iT^i \label{2.9.13}
\end{equation} The $E_{\alpha}, D_i $ and $T^i $ are as defined in (\ref{2.2.4}), (\ref{2.2.2}) and (\ref{2.9.11}) respectively. The right hand side of (\ref{2.9.13}) is a divergence expression. To find the multipliers $\Delta^{\alpha}$, one invokes and solves the determining equations (\ref{2.9.14}) (\cite{bluman2008symmetry}) \begin{equation} \frac{ \delta }{ \delta u^{\alpha}} ( \Delta^{\alpha} E_{\alpha} )=0 \label{2.9.14}
\end{equation}

\subsection{Ibragimov's conservation theorem}
This approach \citep{ibragimov2009practical}, enables one to derive a conservation law for each Lie point symmetry of the partial differential equation system.

 \begin{defn}  Consider the partial differential equation system (\ref{2.2.4}). The adjoint equations of the system (\ref{2.2.4})  are defined by \begin{equation} E_{\alpha}^{*}( x, u , v, \ldots,u_{(\pi)},v_{(\pi)}) \equiv   \frac{ \delta }{ \delta u^{\alpha }} ( v^{\sigma } E_{\sigma})   =0, \label{2.9.15}
	\end{equation} wherein  $ v= ( v^1, \ldots, v^m)$ is the new dependent variable.
>>>>>>> f3e458080f41e5290173f14b8121103bb6a456e0
\end{defn} 
\begin{defn} Formal Lagrangian $ \mathcal{L}$ of the system (\ref{pi-thorderpde}) and its adjoint equations  (\ref{sec2.9.15}) is \citep{ibragimov2007new} \begin{equation} \mathcal{L} = v^{\alpha} \Delta_{\alpha } ( x^i,u^{\alpha},u_{(1)}, \ldots, u_{(\pi)}). \label{sec2.9.16}
\end{equation}
<<<<<<< HEAD
\end{defn}
\begin{thm} \label{conls} Every infinitesimal symmetry $
X $of the system given by (\ref{pi-thorderpde}) leads to  conservation laws \citep{ibragimov2007new} \begin{equation}
	D_iT^i\Big|_{\Delta_{\alpha =0 }} =0,  \label{sec2.9.18}
	\end{equation} where the conserved vector \begin{align}  \begin{aligned}
=======

\begin{thm} \citep{ibragimov2009practical} Every infinitesimal symmetry $
X $of the system (\ref{2.2.4}) leads to a a conservation law for the system ( \ref{2.2.4}) \begin{equation}
	D_iT^i|_{E_{\alpha =0 }} =0  \label{2.9.18}
	\end{equation}. The components of the conserved vector $ T^i$ are given by the formula \begin{align}  \begin{aligned}
>>>>>>> f3e458080f41e5290173f14b8121103bb6a456e0
	T^i = \xi^i \mathcal{L} +  W^{\alpha} \left[  \frac{ \partial \mathcal{L}}{ \partial u^{\alpha}_i} - D_j \left(\frac{ \partial \mathcal{L}}{ \partial u^{\alpha}_{ij}} \right) + D_jD_k \left(\frac{ \partial \mathcal{L}}{ \partial u^{\alpha}_{ijk}} \right)- \ldots    \right] + \\
	D_j ( W^{\alpha} ) \left[ \frac{ \partial \mathcal{L}}{ \partial u^{\alpha}_{ij}} -D_k  \left(\frac{ \partial \mathcal{L}}{ \partial u^{\alpha}_{ijk}} \right)+ \ldots   \right] + 
	D_jD_k ( W^{\alpha} ) \left[ \frac{ \partial \mathcal{L}}{ \partial u^{\alpha}_{ijk}} - \ldots   \right].\end{aligned}
\end{align} 
\end{thm}
<<<<<<< HEAD
\section{Concluding remarks}
This chapter has introduced fundamental concepts of Lie group analysis of PDEs and conservation laws. We have defined local Lie groups, their generators, and infinitesimal transformations. We have also shown how one can obtain prolongations to a generator and prolonged groups. This is very key when solving determining equations for symmetries of a PDE. The ideas of group-invariants and characteristic equations will help to find group-invariant solutions.
Also important to mention is the methods of obtaining conservation laws of PDEs. For this dissertation, we have considered only the method of multipliers and Ibragimov's conservation theorem.

=======
\section{Conclusions}
In this chapter, we presented a succinct introduction to Lie group analysis and conservation laws of partial differential equations and gave some concrete results which will be used throughout this dissertation. We also touched on some algorithms of certain methods used to determine conservation laws. 
>>>>>>> f3e458080f41e5290173f14b8121103bb6a456e0
 % Chapters might go from 2. problem statement, 
                 % through 3. model, to 4. analysis & results
\chapter{Korteweg-de Vries (KdV)  Equation \label{ch3}}
We introduced the concepts of Lie group analysis and conservation laws in Chapter \ref{ch2}, which we will rely on in this chapter. Our aim is to provide an illustrative example of Lie group analysis of a nonlinear partial differential equation (NLPDE). We find  Lie point symmetries of the KdV equation and use them for symmetry reductions and construction of invariant solutions. In addition, we use the multiplier method to get its conservation laws.
\section{Introduction}
A Scottish civil engineer, John Scott Russell is credited with first work on solitons \citep{wazwaz2010partial}. Russell observed that water waves maintained their shape and structure as they traveled in a canal. He went ahead to conduct experiments, which culminated in the discovery of solitons.
In the year 1895, Diederik Korteweg and Gustav De Vries analytically derived the  KdV equation. However,  Joseph  Boussinesq had earlier (1877) introduced the equation in his work on water waves. The KdV equation has applications in aerodynamics, fluid dynamics, and continuum mechanics among other disciplines of mathematical sciences. The equation has been used to describe the dynamics of solitons, ion-acoustic waves in plasmas, surface waves in shallow water, and long internal waves in oceans. The KdV equation also models shock wave formation, turbulence, boundary layer behavior, and mass transport. The simplest form of the KdV equation is  given by \begin{equation}
u_t + a uu_x + u_{xxx} =0.  \label{ew0}\end{equation} The KdV equation combines a quadratic nonlinear dissipative term $uu_x$ which localizes the wave and a linear dispersive term $u_{xxx}$ which spreads  out the wave. The physical meaning  of $u(t,x)$ is the local elevation of the wave surface  at time $t$ and position $x$. With the help of scaling, we associate any equation of the form
\begin{equation}
\alpha u_{t}+\beta u u_{x}+ \gamma u_{xxx}=0,
\end{equation}
to be of “KdV type”.  In this chapter, we use the  form
\begin{equation}\label{kd}
\Delta \equiv  u_{t}+6u u_{x}+ u_{xxx}=0,
\end{equation} where the 6 factor is just conventional and of no great significance. Infact, most commonly used factors are $ \pm 1 , \pm 6.$
\section{Solutions of the KdV Equation (\ref{kd})}
\subsection{Lie point symmetries of \label{secv} (\ref{kd}) }

We use  Lie's method to determine Lie point symmetries of the KdV equation. The  infinitesimal transformations of a local  Lie group with a parameter $\epsilon$ are 
\begin{align}
\begin{aligned}
\bar{t}=t+  \tau(t,x,u)\epsilon,\,\,\bar{x}=x+  \xi(t,x,u)\epsilon, \,\,
\bar{u}=u+  \eta(t,x,u)\epsilon.
\label{infintlil} 
\end{aligned}
\end{align}
The vector field
\begin{eqnarray}\label{kdv1}
X =\tau (t,x,u) \frac{\partial}{\partial t} + \xi (t,x, u) \frac{\partial}{\partial x} 
+\eta (t,x,u) \frac{\partial}{\partial u} , 
\end{eqnarray}
is a Lie point symmetry of  (\ref{kd}) provided that
\begin{eqnarray}
\mbox{X}^{[3]}  \Delta \Big|_{\Delta=0} =0 \label{eq},
\end{eqnarray} where
\begin{eqnarray}
	\mbox{X}^{[3]} =X +\zeta_{1}\frac{\partial }{\partial u_{t}}+\zeta_{2}\frac{\partial }{\partial u_{x}}+\zeta_{222}\frac{\partial }{\partial u_{xxx}},
\end{eqnarray} is the third prolongation of the Lie point symmetry $X$ as defined in  (\ref{sec2.4.3}) of Chapter \ref{ch2} and
\begin{eqnarray}
	\zeta_1 &=& D_t (\eta) - u_t D_t (\tau) - u_x D_t (\xi), \\
	\zeta_2 &=& D_x (\eta) - u_t D_x (\tau) - u_x D_x (\xi),\\
	\zeta_{22} &=& D_x(\zeta_2) - u_{tx}D_x(\tau) - u_{xx}D_x(\xi),\\
	\zeta_{222} &=& D_x(\zeta_{22}) - u_{txx}D_x(\tau) - u_{xxx}D_x(\xi),
\end{eqnarray} as defined in  (\ref{sec2.4.2}), and
\begin{eqnarray}
\label{f1} D_t &=& \frac{\partial}{\partial t} + u_t\frac{\partial}{\partial u}
+ u_{tx} \frac{\partial}{\partial u_x} + u_{tt}\frac{\partial}{\partial u_t} + \cdots \, , \\
\label{f2} D_x &=& \frac{\partial}{\partial x} + u_x \frac{\partial}{\partial u} + u_{xx}\frac{\partial}
{\partial u_x} + u_{tx} \frac{\partial}{\partial u_{t}} + \ldots .
\end{eqnarray}
Applying the definitions of $D_t$ and $D_x$ given in (\ref{f1}) and (\ref{f2}), we obtain the expanded form of the $\zeta_s$  as \vspace{0.005cm}
\begin{align}  \begin{aligned}
\zeta_1& = \eta_t + u_t\eta_u - u_t\tau_t - u_t^2\tau_u - u_x\xi_t - u_tu_x\xi_u,\\ 
\zeta_2& = \eta_x + u_x \eta_u - u_t \tau_x - u_t u_x \tau_u - u_x\xi_x - u_x^2\xi_u,\\
\zeta_{22}& = \,  \eta_{xx} + 2u_x \eta_{xu} + u_{xx} \eta_u + u_x^2 \eta_{uu} - 2u_{xx} \xi_x
-u_x \xi_{xx} - 2u_x^2 \xi_{xu} - 3u_xu_{xx} \xi_u \\
&
- u_x^3 \xi_{uu}- 2u_{tx} \tau_x - u_t \tau_{xx}  - 2u_tu_x \tau_{xu} - (u_tu_{xx} + 2u_xu_{tx}) \tau_u - u_t u_x^2 \tau_{uu},\\
\zeta_{222}&= \eta_{xxx}
-3u_{t}u_xu_{t}u_{xx}\tau_{uu}+u_{x}^3\eta_{uuu}+3u_{x}^2\eta_{uux}-u_t\tau_{xxx}-3u^2_x\xi_{uux}-3u_{x}^3 \xi_{uux} \\
&
-u_{x}^4\xi_{uuu}-3u_{xx}^2 \xi_{u}+3u_{xx} \eta_{ux} +u_{xxx}\eta_u-3u_{txx}\tau_{x} -3u_{tx}\tau_{xx}-3u_{xxx} \xi_{x}\\
&
-3u_{xx}\xi_{xx}-3u_t u_{x}\tau_{uxx} -3u_t u_{xx} \tau_{ux}-6u_{tx} u_x \tau_{ux}+3u_{xx} u_x\eta_{uu}-3u_{tx}u_{x}^2 \tau_{uu} \\
&
+3u_{x}\eta_{uxx}-u_{x}\xi_{xxx}-3u_{t} u_{x}^2\tau_{uux}-u_{t} u_{xxx}\tau_u -3u_{txx} u_x\tau_u-u_{t} u_{x}^3\tau_{uuu}\\
&
-4u_{xxx} u_x\xi_u-9 u_{x} u_{xx}\xi_{ux}-3u_{tx} u_{xx}\tau_u-6u_{x}^2 u_{xx}\xi_{uu}.
\end{aligned}
\end{align}
From (\ref{eq}), we get
\begin{align}
\Big[\tau\frac{\partial}{\partial t} + \xi\frac{\partial}{\partial x}
+ \eta\frac{\partial}{\partial u} + \zeta_1\frac{\partial}{\partial u_t} + \zeta_2
\frac{\partial}{\partial u_x} + \zeta_{222}\frac{\partial}{\partial u_{xxx}}\Big]\Big(u_{t}+6 uu_x+ u_{xxx}\Big)\Big|_{ (\label{kd1})} = 0.
\end{align}
After expanding   (\ref{kd1}), we obtain
\begin{align} \label{6u}
(6u_x \eta +\zeta_1 +6u \zeta_{2}+ \zeta_{222})\Big|_{u_{xxx}=-u_{t}-6 uu_x} = 0.
\end{align}
Substitutions of values of $\zeta_1$, $\zeta_2$ and $\zeta_{222}$ in   (\ref{6u}) yield
\begin{align}\begin{aligned} \label{6x}
& 
6u_x\eta+ [\eta_t +u_t \eta_u -u_t\tau_u-u^2_t \tau_u-u_x \xi_t-u_tu_x \xi_u ] + 6u[\eta_x +  u_x \eta_u- u_t\tau_x- u_t u_x\tau_u\\
&
- u_x\xi_x - u^2_x\xi_u]+[\eta_{xxx}
-3u_{t}u_xu_{t}u_{xx}\tau_{uu}+u_{x}^3\eta_{uuu}+3u_{x}^2\eta_{uux}-u_t\tau_{xxx}-3u^2_x\xi_{uux} \\
&
-3u_{x}^3 \xi_{uux}-u_{x}^4\xi_{uuu}-3u_{xx}^2 \xi_{u}+3u_{xx} \eta_{ux} +u_{xxx}\eta_u-3u_{txx}\tau_{x} -3u_{tx}\tau_{xx}-3u_{xxx} \xi_{x}\\
&
-3u_{xx}\xi_{xx}-3u_t u_{x}\tau_{uxx}-3u_t u_{xx} \tau_{ux}-6u_{tx} u_x \tau_{ux}+3u_{xx} u_x\eta_{uu}-3u_{tx}u_{x}^2 \tau_{uu} \\
&
-u_{x}\xi_{xxx}-3u_{t} u_{x}^2\tau_{uux}-u_{t} u_{xxx}\tau_u -3u_{txx} u_x\tau_u-u_{t} u_{x}^3\tau_{uuu}-4u_{xxx} u_x\xi_u-9 u_{x} u_{xx}\xi_{ux}\\
&
-3u_{tx} u_{xx}\tau_u-6u_{x}^2 u_{xx}\xi_{uu}] \Big|_{ u_{xxx}=-u_{t}-6 uu_x}  =0.\end{aligned}
\end{align}
Replacing  $u_{xxx}$ by $-u_{t}-6 uu_x $  in (\ref{6x}), we obtain
\begin{align}\begin{aligned}
&6u_x\eta+ [\eta_t +u_t \eta_u -u_t\tau_u-u^2_t \tau_u-u_x \xi_t-u_tu_x \xi_u ] + 6u[\eta_x +  u_x \eta_u- u_t\tau_x- u_t u_x\tau_u\\
&
- u_x\xi_x - u^2_x\xi_u]+[\eta_{xxx}
-3u_{t}u_xu_{t}u_{xx}\tau_{uu}+u_{x}^3\eta_{uuu}+3u_{x}^2\eta_{uux}-u_t\tau_{xxx}-3u^2_x\xi_{uux} \\
&
-3u_{x}^3 \xi_{uux}-u_{x}^4\xi_{uuu}-3u_{xx}^2 \xi_{u}+3u_{xx} \eta_{ux} +(-u_{t}-6 uu_x)\eta_u-3u_{txx}\tau_{x} -3u_{tx}\tau_{xx}\\
&
-3(-u_{t}-6 uu_x) \xi_{x}-3u_{xx}\xi_{xx}-3u_t u_{x}\tau_{uxx} -3u_t u_{xx} \tau_{ux}-6u_{tx} u_x \tau_{ux}+3u_{xx} u_x\eta_{uu}\\ 
&
-3u_{tx}u_{x}^2 \tau_{uu}+3u_{x}\eta_{uxx}-u_{x}\xi_{xxx}-3u_{t} u_{x}^2\tau_{uux}-u_{t} (-u_{t}-6 uu_x)\tau_u -3u_{txx} u_x\tau_u \\
&-u_{t} u_{x}^3\tau_{uuu}-4(-u_{t}-6 uu_x) u_x\xi_u-9 u_{x} u_{xx}\xi_{ux}-3u_{tx} u_{xx}\tau_u
-6u_{x}^2 u_{xx}\xi_{uu}]=0. \label{6t}\end{aligned}
\end{align}
 Splitting  (\ref{6t}) on  derivatives of $u$ gives an overdetermined system of eight PDEs, namely
\begin{eqnarray}
\label{1}u_{txx} &:& \tau_x=0, \\
\label{2}u_{tx} u_{xx} &:& \tau _u=0, \\
\label{3}u_{xx}^2& :&\xi _u=0, \\
\label{7}u_x u_{xx}& : &\eta_{uu}=0,\\
\label{4}u_{xx} &: &\eta_{ux}-\xi_{xx} =0, \\
\label{5}u_{t}& : &3\xi_x -\tau_{t} =0, \\
\label{6}u_{x} &:& 6\eta + 12 u \xi_x - \xi_{t}+3 \eta_{uxx}-\xi_{xxx} =0,\\
\label{8}\mbox{Rest} &:& \eta_{xxx}+6u\eta_{x}+\eta_{t}=0.
\end{eqnarray}
From (\ref{1}) and (\ref{2}), we have \begin{align} \label{pau}	\tau = \tau(t), \end{align} where  arbitrary function $\tau$ depends on $t$.
Substituting the expression of $\tau$ given by (\ref{pau}) into (\ref{5}), and integrating thereafter with respect to $x$ while taking note of (\ref{3}), we have 
  \begin{equation} \xi = \frac{\tau_t(t) }{3}x + a(t),\label{sec3.2.10}
\end{equation} for some arbitrary function $a$ of $t$. 
Using the expression  for $\xi$ given by (\ref{sec3.2.10}) in  (\ref{6}), we have 
\begin{align}\eta = \frac{-2\tau_{t}(t) u}{3}+  \frac{\tau_{tt}(t) }{18}x +  \frac{a_t(t)}{6}. \label{bau}
\end{align}
The expressions for  $\xi$ and $\eta$ as given by  (\ref{sec3.2.10}) and  (\ref{bau}) respectively, 
 admit (\ref{4}).  Substituting the expression of $\eta$ given by (\ref{bau}) in (\ref{8}) yields 
\begin{align} \label{las}-\frac{\tau_{tt}(t) }{3}u + \frac{\tau_{ttt}(t) }{18}x + \frac{ a_{tt}(t)}{6} =0. 
\end{align} 
 If we split  (\ref{las}) on powers of $u$ and $x$, we obtain
 \begin{eqnarray}
\label{pau1} u&:& \tau_{tt}(t) =0,\\
\label{pau2}  x &:& \tau_{ttt}(t) =0,\\
\label{pau3} \mbox{rest}  &:& a_{tt}(t)=0.
 \end{eqnarray}
We observe that  (\ref{pau1}) satisfies  (\ref{pau2}). Integration of  (\ref{pau1}) and (\ref{pau3})  twice with respect to $t$ gives \begin{align}
 \tau &= 3C_1t + C_2,\\
 a(t)&= 6C_3 t + C_4.
 \end{align}
Finally, appropriate substitutions give 
\begin{align}
\label{111}  \tau =& 3C_1t + C_2, \\
\label{121}  \xi = &C_1x + 6C_3t + C_4, \\
\label{131} \eta = &- 2C_1u+C_3. 
\end{align}
The  calculations  prove that the KdV Equation (\ref{kd}) admits  a four-dimensional Lie algebra generated by
\begin{align}
 \label{sec3.2.22} X_1 = &\frac{ \partial}{ \partial x},\\
\label{sec3.2.28} X_2 =& \frac{ \partial}{ \partial t},\\
 \label{sec3.2.32} X_3 =& 6t\frac{ \partial}{ \partial x } + \frac{ \partial }{ \partial u},\\
\label{sec3.2.35} X_4 =&3t \frac{ \partial }{ \partial t}+ x\frac{ \partial}{ \partial x}  - 2 u \frac{ \partial }{ \partial u}.
\end{align}
 \begin{rem}
The first two symmetries represent space and time translations respectively  while the third represents Galilean boost and the fourth represents scaling symmetry.
 \end{rem}
 \subsection{Commutator table} The set of all solutions to the determining Equation (\ref{eq}) forms a Lie algebra.
 In this subsection , we evaluate the commutation relations for the symmetry generators.  We use the definition of Lie bracket in  (\ref{sec2.7.2}). For example, we have that the commutator \begin{align}
 \left[X_1, X_2\right] = \left( \frac{ \partial}{ \partial x} \frac{ \partial}{ \partial t}\right)- \left( \frac{ \partial}{ \partial t} \frac{ \partial}{ \partial x} \right) =0.
 \end{align}
  \begin{rem}
  The remaining commutation  relations  are obtained similarly as displayed in Table   \ref{Tab:Tcr}. 
 \begin{center}
 	\begin{tabular}{|c|c|c|c|c|} 
 		\hline 
 		$[X_i,X_j]$	& $ X_1$ & $ X_2$ & $ X_3$ & $ X_4$ \\ 
 		\hline 
 		$ X_1$ & 0 & 0 & 0 & $ X_1$ \\ 
 		\hline 
 		$ X_2$ & 0 & 0 & 6 $ X_1$ & 3$ X_2$ \\ 
 		\hline 
 		$ X_3$ & 0 & -6$ X_1$ & 0 & -2 $ X_3$  \\ 
 		\hline 
 		$ X_4$ & -$X_1$ & -3$ X_2$ & 2$ X_3$ & 0 \\ 
 		\hline 
 	\end{tabular}  \captionof{table}{\label{Tab:Tcr}A commutator table for the Lie algebra generated by the symmetries of the KdV Equation (\ref{kd}).}
 \end{center}
  \end{rem}
\subsection{Local Lie groups of transformations}
The corresponding one-parameter groups of transformations are determined by solving the Lie equations as given in Theorem (\ref{lieeqns}). Let $T_{\epsilon_i}$  be the group of transformations  for each $X_i, i =1,2,3,4. $   To obtain  $T_{\epsilon_1}$ from the infinitesimal generator $X_1$ (\ref{sec3.2.22}), one integrates the Lie equations
\begin{align} \begin{aligned}
\frac{\mathrm{d} \bar{t} }{ \mathrm{d} \epsilon_1 } = 0,\quad \bar{t}\Big |_{ \epsilon_1 =0} =t, \,\,
\frac{ \mathrm{d} \bar{x}}{ \mathrm{d}\epsilon_1} =1,\,\, \bar{x}\Big |_{ \epsilon_1 =0} =x, \,\,
\frac{ \mathrm{d} \bar{u}}{ \mathrm{d}\epsilon_1} = 0,\,\, \bar{u}\Big |_{ \epsilon_1 =0} =u.
\label{lieeq3} \end{aligned}
\end{align}
Solving the  system (\ref{lieeq3}), one obtains
\begin{align} T_{\epsilon_1} : \,\,\bar{t}=t,\,\,  \bar{x}=x + \epsilon_1,\,\, \bar{u} =u. \label{sec3.1.14}
\end{align}
The other three local Lie  groups are obtained similarly and are given by
\begin{align} \begin{aligned}
T_{\epsilon_2} &: \,\,  \bar{t} = t+\epsilon_2,\,\,  \bar{x}=x,\,\,  \bar{u}=u, \\
T_{\epsilon_3} &:\,\,  \bar{t}=t, \,\, \bar{x}=x + 6\epsilon_3t,\,\,  \bar{u}=u+ \epsilon_3,\\
T_{\epsilon_4} &:\,\,  \bar{t}=t e^{3 \epsilon_4}, \,\, \bar{x}=x e^{\epsilon_4},\,\, \bar{u} =ue^{-2\epsilon_4}.
\label{sec3.1.15} 
\end{aligned}
\end{align}
\begin{rem}
 In Lie group analysis, one can obtain solutions of PDEs either by transforming known solutions by the one-parameter groups or constructing group-invariant invariant solutions. 
\end{rem}
\subsection{Transformation of known solutions}
We show how to transform known solutions.
By the criterion of invariance, if $ \bar{u} = \Gamma( \bar{t}, \bar{x})$ admits (\ref{kd}), then so does \begin{equation}
\phi(t,x,u,\epsilon) = \Gamma ( \varphi_1(t,x,u,\epsilon) \varphi_2(t,x,u,\epsilon) ). \label{sec3.1.16}
\end{equation}
Using the Lie groups obtained in (\ref{sec3.1.14}) and (\ref{sec3.1.15}), one gets  the transformed solutions 
\begin{align} \begin{aligned} 
T_{\epsilon_1}:& u^{(1)} = \Gamma(t,x + \epsilon_1),\\
T_{\epsilon_2}:& u^{(2)} = \Gamma(t + \epsilon_2,x),\\
T_{\epsilon_3}: &u^{(3)} = \Gamma(t,x + 6 t\epsilon_3) - \epsilon_3,\\
T_{\epsilon_4}: &u^{(4)} =e^{2 \epsilon_4} \Gamma(te^{3 \epsilon_4},x e^{\epsilon_4}).\\
\label{sec3.1.17} \end{aligned}
\end{align}
 \subsection{Group-invariant solutions of (\ref{kd})}
 We  compute the group-invariant solutions of the KdV Equation (\ref{kd}) under all its Lie point symmetries.
 \begin{enumerate}[(i)]
 \item \textbf{Space Translation-Invariant Solutions.} \newline
 Consider the space translation operator 
$ X_1 = \partial/ \partial x.$  The  characteristic equations defined in (\ref{sec2.6.4}) of Chapter \ref{ch2} for  the operator (\ref{sec3.2.22}) are \begin{equation} \frac{ \mathrm{d} t }{0 } = \frac{ \mathrm{d} x }{1 } = \frac{ \mathrm{d} u }{0 }, \label{sec3.2.23}
 \end{equation} which gives two invariants $ J_1 =t$ and $ J_2 = u$. Therefore, $J_2 = \psi(J_1) $  is the group-invariant solution for some arbitrary function $ \psi$. Substitution of $u=\psi(t)$ into (\ref{kd}) yields $
 \psi'(t) =0,$  whose solution is $ \psi(t) = C_1,$  for  some arbitrary constant $C_1$. Hence the space translation-invariant solution of (\ref{kd})  is \begin{equation}
u(t,x) = C_1. \label{sec3.2.27}
 \end{equation}
\item \textbf{Time translation-invariant (Stationary) solutions.} \newline
 Consider the time translation operator $X_2 =  \partial / \partial t .$  The Lagrangian system associated with the operator (\ref{sec3.2.28}) yields the invariants  $J_1 = x$ and $ J_2 = u$. 
Thus, $ u = \psi(x)$ is the group-invariant solution.
Substituting of $u= \psi(x)$ into (\ref{kd}) yields \begin{equation} 6 \psi'(x) \psi(x) + \psi'''(x)=0. \label{ssol}
\end{equation} All the stationary solutions have the form $ u =  \psi(x)$
 for some arbitrary function $ \psi$  satisfying (\ref{ssol}) or 
 \begin{equation}
(\psi')^2 = -2 \psi^3 + 2k \psi +l, \label{psi}
\end{equation} obtained from (\ref{ssol}) by two integrations, for which $ k $ and $l$ are constants of integration.\newline The general invariant solution takes the form \citep*{ibragimov1995crc} \begin{equation}
u = -2 \Phi(x),  \label{ssta}
\end{equation} where $\Phi(x)$ is the Weierstrass elliptic function satisfying \begin{equation}  \Phi(x)'^2 = 4  \Phi(x)^3 -g_2 \Phi(x)-g_3. \label{secwes}
\end{equation} For some real roots $r_1, r_2 ,r_3$, of the cubic polynomial on the right-hand side of (\ref{psi}), the solutions (\ref{ssta}) could be rewritten in the following forms: \citep{ibragimov1995crc}
\begin{enumerate}[(a)]
	\item If $r_1 <r_2 <r_3$, then $ u = u(x)$, is a limited function and \begin{equation}\label{sba} u = \frac{2a}{s^2} \mathrm{dn^2} \left( \sqrt{ \frac{a}{s^2}}x,s\right)+ \gamma,
	\end{equation} is a cnoidal wave where $\mathrm{dn^2}$(x,s) is the Jacobian elliptic function with modulus \newline$ s = \sqrt{ \frac{r_3-r_2}{  r_3-r_1} },$ where $
	 a=\frac{r_3-r_2}{2}$ is the amplitude of a wave, and $ \gamma=r_1$ .
	\item If $ r_1 = r_2 < r_3, $ then $ u \to r_1, u', u'' \to 0$ when $ |x| \to \infty$   \begin{equation} u = (r_3-r_1) \sech^2 \left( \sqrt{ \frac{r_3-r_1}{2}}\right)+ r_1,
	\end{equation} is a  solitary wave.
	\item If $ r_1 = r_2 =r_3$, then 
	\begin{equation} u = -\frac{2}{ (x-c)^2},\quad  x\neq c.
	\end{equation}
\end{enumerate}
\item \textbf{Galilean-invariant solutions.} \newline
 Consider the Galilean boost operator $ X_3 = 6t  \partial / \partial x + \partial / \partial u. $
  The characteristic equations associated to the operator (\ref{sec3.2.32}) yield two invariants $ J_1 = t$ and  $J_2 = -u+x/(6t)$. As a result, the group-invariant solution of (\ref{kd}) for this case is $ J_2 = \phi(J_1)$, for $\phi$ an arbitrary function. That is, \begin{equation} u(t,x) =-\phi(t)+ \frac{x}{6t}, \quad t\neq 0. \label{sec3.2.33}
\end{equation}
Substitution of the expression for $u$ from  (\ref{sec3.2.33}) into  (\ref{kd}) yields a first order ordinary DE $ \phi'(t) +  \phi(t)/t=0,$ whose general solution is  $\phi(t) =  \delta/ t $ for  an arbitrary constant $\delta$. Hence, we have a group-invariant solution for $X_3$ as  \begin{equation} u(t,x) =   \frac{ x+ \mathcal{C}}{6t} , \label{sec3.2.34}
\end{equation}  where $ \mathcal{C}=-6\delta$ and $t \neq 0.$
\item \textbf{Scale-invariant solutions}.\newline
Last but not least, we consider the scaling operator $
X_4 = x \partial /\partial x + 3t \partial / \partial t -2u  \partial/ \partial u,$  whose associated Lagrangian equations yield two invariants, $J_1 = x^3/t$ and $J_2 = ux^2$. Thus, 
 \begin{equation} u = x^{-2}  \varphi(\lambda) , \quad \lambda=\frac{x^3}{t},  \label{ssol1} \quad  t \neq 0,
\end{equation} is the group-invariant solution under $X_4$ where  $\varphi$  satisfies  
\begin{eqnarray}\label{17}
27\lambda^3 \varphi'''	+18\lambda \varphi \varphi'+(24 - \lambda) \lambda \varphi' -12(2 + \varphi)\varphi =0.
\end{eqnarray}
\end{enumerate}
\subsection{Travelling wave solution} 
To obtain travelling wave solution, in particular a soliton solution of the KdV equation, we consider a linear combination of the space and time translation symmetries, namely,  $ X = c {\partial }/{\partial x} +  {\partial }/{\partial t}$ for some arbitrary constant $c$ considered as the velocity of the wave.
The characteristic equations give  two invariants, $ J_1 = u$ and $ J_2 = x -ct$.  So  $ u(t,x) = \varphi(x-ct)$, for some arbitrary function  $\varphi$,  is the invariant solution.
Substitution of $u$ into (\ref{kd}) yields a third-order ordinary DE 
\begin{equation} 
-c\varphi'(\xi)+ 6 \varphi (\xi)\varphi' (\xi)+ \varphi'''(\xi) = 0, \,\,\, \xi = x-ct. \label{sec3.2.40}
\end{equation}
Integrating  (\ref{sec3.2.40}) with respect to $\xi$ yields \begin{equation}  
-c\varphi (\xi)+ 3 \varphi^2 (\xi)+ \varphi''(\xi) = 0 \label{sec3.2.41}, 
\end{equation} 
where we have chosen  $0$ as the constant of integration. The second integration is done after multiplying  (\ref{sec3.2.41}) by $ 2 \varphi'$ and we get 
\begin{align} 
(\varphi')^2 = c \varphi^2- 2 \varphi^3 \quad  \text{ or}  \quad \frac{ \mathrm{d} \varphi}{ \sqrt{ c \varphi^2- 2 \varphi^3  } } = \mathrm{d} \xi. \label{oti}
\end{align} 
In the second integration, again we took the constant of integration to be $0$. By the change of variable $ \varphi =({c}/{2}) \sech^2( \eta)$ and integration of (\ref{oti}), we get the one-soliton solution 
\begin{equation}
u(x,t) = \frac{c}{2} \sech^2\left( \frac{ \sqrt{c}}{2}(x-ct)\right).
\end{equation}
\section{Conservation laws of the KdV Equation (\ref{kd})}
Computation of conservation laws for KdV equation  (\ref{kd}) is done by  the using the  method of multipliers \citep{olver2000applications}. 
\subsection{Conservation Laws of the KdV Equation (\ref{kd}) using multipliers \label{consl}}
Using  the Euler-Lagrange operator defined  in (\ref{sec2.9.1}), we   search for a zeroth order multiplier \newline $\Lambda(t,x,u)= \Lambda $. The  determining equation for computing $ \Lambda$ is \begin{equation}
\frac{ \delta}{ \delta u} \left[ \Lambda \{  u_t + 6 u u_x + u_{xxx}  \}\right] =0. \label{d}
\end{equation}
Expansion of  (\ref{d}) yields 
\begin{equation}
\Lambda_u ( u_t + 6 u u_x + u_{xxx} ) + 6 u_x \Lambda-D_t (\Lambda)-6D_x(\Lambda u)-D_x^3(\Lambda)=0. \label{sec3.3.2}
\end{equation}
Invoking the total derivatives defined in (\ref{f1}) and (\ref{f2}) on  (\ref{sec3.3.2}), we get
\begin{align}
\Lambda_t + 6 u \Lambda_x + \Lambda_{xxx}+ 3\Lambda_{xxu} u_x +  3 \Lambda_{xuu}   u_x^2 
+ \Lambda_{uuu} u_x^3 + 3 \Lambda_{xu} u_{xx} +  3 \Lambda_{uu}  u_x u_{xx}=0. \label{sec3.3.3} \end{align}
Splitting (\ref{sec3.3.3}) on derivatives of $u$ produces a simplified system of 3 PDEs, namely,
\begin{align}
u_{xx} &: \Lambda_{xu}=0, \label{sec3.3.7} \\
u_x u_{xx}&: \Lambda_{uu}=0, \label{sec3.3.8}\\
rest&:\Lambda_t + 6 u \Lambda_x + \Lambda_{xxx}=0. \label{sec3.3.9}
\end{align}
Integrating (\ref{sec3.3.8})  with respect to $u$ twice gives \begin{equation} \Lambda =  \mathfrak{a}(t,x)u + \mathfrak{b}(t,x),  \label{sec3.3.10} \end{equation}
 for some arbitrary functions $\mathfrak{a}$  and  $\mathfrak{b}$ of  $t$ and $x$.
 Substitution of the value of $\Lambda$ from  (\ref{sec3.3.10})  into  (\ref{sec3.3.7}) yields, \begin{equation}  \mathfrak{a}_x(t,x) =0, \label{aq}
	 \end{equation} from which we get that   $\mathfrak{a}= \mathfrak{a}(t)$, and consequently \begin{equation}
	 \Lambda =  \mathfrak{a}(t)u + \mathfrak{b}(t,x). \label{sec3.3.11}
	 \end{equation}
	 If we substitute  the expression for $\Lambda$ given in   (\ref{sec3.3.11}) into  (\ref{sec3.3.9}), we obtain \begin{equation} \mathfrak{a}_t(t) u + \mathfrak{b}_t(t,x) + 6 u\mathfrak{b}_x(t,x)+ \mathfrak{b}_{xxx}(t,x) =0. \label{sec3.3.12}
\end{equation} 
 Splitting  (\ref{sec3.3.12}) on powers of $u$, we find	
\begin{eqnarray}
\label{sec3.3.13}u &:& \mathfrak{a}_t(t)+ 6 \mathfrak{b}_x(t,x)=0, \\
\label{sec3.3.14}u^0 &:&  \mathfrak{b}_t(t,x)+ \mathfrak{b}_{xxx}(t,x) =0.
\end{eqnarray}	
Equation (\ref{sec3.3.13})  implies that 	\begin{equation}
\mathfrak{b}_x(t,x)= -\frac{ \mathfrak{a}_t(t) }{6} , \label{sec3.3.15}
\end{equation} from which $ \mathfrak{b}_{xxx}(t,x)=0.$ Thus, by  (\ref{sec3.3.14}), we deduce that $ \mathfrak{b}_t(t,x)=0$, meaning  that 
$\mathfrak{b}=\mathfrak{b}(x).$
\ Upon integration of (\ref{sec3.3.15}) with respect to $x$, we have that 
\begin{equation}
	\mathfrak{b}(x)=-\frac{ \mathfrak{a}_t(t) }{6} x+ C_1,   \label{sec3.3.17}
\end{equation} for some arbitrary constant $C_1$. Since
  $\mathfrak{a}=\mathfrak{a}(t)$ as shown by  (\ref{aq}), we must have that  $\mathfrak{a}_t(t)=C_2 $ for some arbitrary constant $C_2$.
 This gives $ \mathfrak{a}(t)=C_2 t + C_3$ and $	\mathfrak{b}(x)= C_1-C_2x/6  $  for an  arbitrary constant $C_3$.
Hence  we have, \begin{equation} \Lambda = C_1 +C_2\left( tu-\frac{x }{6} \right) + C_3 u,
\end{equation} which means we have  three non-trivial conservation law multipliers, namely,\begin{align}
 \Lambda_1 = 1, \quad 
\Lambda_2= tu-\frac{1 }{6}x \quad \text{and} \quad \Lambda_3 = u.	
\end{align} 
\begin{rem} Recall from (\ref{sec2.9.13}) that
a multiplier  $\Lambda $ for  (\ref{kd}) has the property that   for  the density $ T^t = T^t(t,x,u,u_x) \quad \text{and flux}  \quad  T^x = T^x(t,x,u,u_x,u_{xx}),$  \begin{equation} \Lambda \left( u_t + 6uu_x + u_{xxx}\right) = D_t T^t + D_x T^x \label{sec3.3.20}.
\end{equation} 
\end{rem} We look for  conserved vectors corresponding to each of the three multipliers obtained.
\begin{enumerate}[(a)]
\item \textbf{Conserved vectors for the multiplier  $\Lambda_1= 1.$}
	Expansion of  (\ref{sec3.3.20})  gives,
	\begin{equation}   u_t + 6uu_x + u_{xxx} = T_t^t + u_t T^t_u + u_{tx} T^t_{u_x} + T^x_x+ u_x T^x_u + u_{xx} T^x_{u_x}  + u_{xxx} T^x_{u_{xx}}.  
	\label{j2}
	\end{equation}
	Splitting  (\ref{j2})  on third derivatives of $u$ yields;
	\begin{eqnarray}
	\label{j3} u_{xxx} &:&T^x_{u_{xx}}=1, \\
	\label{j5} \text{Rest} &:& u_{t}+6 u u_{x}= T^t_t+ u_tT^t_u+u_{tx}T^t_{u_x}+T^x_x+u_xT^x_u  +u_{xx}T^x_{u_x}.
	\end{eqnarray}	
	By integrating (\ref{j3})  with respect to $u_{xx}$, we deduce that
	\begin{align}
	T^x=u_{xx}+A(t,x,u,u_x), \label{k4}
	\end{align}  for some arbitrary function  $A$. Substituting the expression of $T^x$ from  (\ref{k4}) into  (\ref{j5}), we have \begin{align}
	u_{t}+6 u u_{x}= T^t_t+u_{t} T^t_u +u_{tx}T^t_{u_x}+A_x+u_xA_u  +u_{xx}A_{u_x}. \label{mbilia}
	\end{align}
Equation (\ref{mbilia})	 splits on second derivatives of $u$ to give;
	\begin{eqnarray}
	\label{fe61} u_{tx} &:&T^t_{u_{x}}=0, \\
	\label{fe31} u_{xx} &:&A_{u_{x}}=0, \\
	\label{fe51} \text{Rest} &:&u_t+ 6 uu_x=T^t_t+u_tT^t_u +A_x +u_{x}A_u .
	\end{eqnarray}
Integrating  (\ref{fe61}) and (\ref{fe31}) with respect to $u_x$	manifests that \begin{align} \label{z1}T^t&=B(t,x,u),\\
\label{z2} A& =A(t,x,u),
\end{align} where $B$ is an arbitrary function of its arguments. Substituting the expressions of $T^t$ and $A$ from (\ref{z1})  and (\ref{z2}) respectively, into  (\ref{fe51}) gives \begin{align}\label{jkay1}
	u_t+ 6 uu_x=B_t+u_tB_u +A_x +u_{x}A_u ,
	\end{align}
	which splits on  the derivatives of $u$ to yield
	\begin{eqnarray}
	\label{fe71} u_{t} &:&B_u=1 , \\
	\label{f771} u_{x} &:&A_u= 6 u, \\
	\label{fe81} \text{Rest} &:& A_x+B_t =0.
	\end{eqnarray} 
By integrating   (\ref{fe71}) and (\ref{f771}) with respect to $u$, we find that \begin{align}
 B&=u+C(t,x),\\
 	A&=3 u^2+D(t,x)
	\end{align}  where $C$ and $D$ are arbitrary functions of their arguments.	Substitution of the expressions of $A$ and $B$ into  (\ref{fe81}) shows  that $D_x+C_t=0$. Since $C$ and $D$ contribute to the trivial part of the  conservation law, we  take $C=D=0$. We get the conserved vector for the multiplier $ \Lambda_1= 1$ as \begin{align}
	 T_1^t= u,\quad  T_1^x=3u^2 + u_{xx}.
	 \end{align}	
\begin{rem}
	Similarly for the multipliers $ \Lambda_{2}$ and $ \Lambda_{3}$, we get the following conserved vectors.
\end{rem}
	\item \textbf{Conserved vector for the multiplier} 	$\Lambda_2 =  tu-x /6$ is given by
 \begin{align}
 T_2^t=  \frac{1 }{2}tu^2-\frac{1}{6}xu, \quad 
 T_2^x=2 u^3t-\frac{1}{2} xu^2+ \frac{1}{6}u_x-\frac{1}{2}tu_x^2+ \left(tu-\frac{x }{6} \right)  u_{xx}.
 \end{align}
 
	\item \textbf{Conserved vector  for the \label{sec2} multiplier $\Lambda_3=  u$} is given by
 \begin{align}
T_3^t=\frac{u^2}{2},\quad 
T_3^x= 2 u^3+ u u_{xx}-\frac{1}{2} u_{x}^2.
\end{align}		
\end{enumerate}
\begin{rem}
It can be verified that\begin{align}D_t T^t_i + D_xT^x_i \Big|_{\Delta =0}=0, \label{dan}
\end{align} for $i =1,2,3$. \end{rem}
\begin{rem} The expressions in (\ref{dan}) are conservation laws for the mass, energy and  momentum  for  the KdV Equation (\ref{kd}).
\end{rem}
\begin{rem} The presence of multiplier   $\Lambda_1=1$ shows that the KdV 
	Equation (\ref{kd}) is itself a conservation law.
\end{rem}
\section{Concluding remarks}
In this chapter, we gave an illustrative example of Lie group analysis of a nonlinear PDE, namely the KdV Equation (\ref{kd}). We achieved that by first constructing time and space translations, Galilean boosts, and scaling symmetry. These symmetries have been used to reduce the KdV equation, a process which led to group-invariant solutions, including a soliton solution of (\ref{kd}). All the solutions found describe different states of the system. Finally, conservation laws for mass, momentum, and energy were obtained using the multiplier approach.    That means those quantities are invariant in the evolution of the KdV system. The approach illustrated in this chapter will be very important for studying the nonlinear system of coupled KdV equations in Chapter \ref{ch4}. % You do not need to ha\usepackage[maxlevel=3]{csquotes}ve exactly 4 chapters.
                 % It is probably a good minimum, with 5 chapters 
                 % average,  and 7 chapters might be a maximum.
\chapter{Nonlinear Coupled  Korteweg-de Vries (KdV) Equations \label{ch4}}
We have demonstrated in Chapter \ref{ch3}, how to find Lie point symmetries, symmetry reductions, and invariant solutions for a nonlinear PDE. As an illustrative example,
we studied the nonlinear Korteweg-de Vries equation (\ref{kd}), for which we found invariant solutions which included a one-soliton solution. Conservation laws for mass, momentum, and energy were also constructed by using multipliers. This chapter will focus on Lie group analysis a nonlinear coupled system of Korteweg-de Vries equations. Our aim is to perform a Lie group analysis of this system. This we achieve by calculating Lie point symmetries, performing symmetry reductions, and constructing invariant solutions for the nonlinear coupled KdV system. Thereafter,  we find conservation laws by use of both the method of multipliers and a theorem due to Ibragimov.
\section{Introduction}
Consider the Kortweg-de Vries equation  \begin{equation} q_t + \alpha q q_x + \beta q_{xxx} =0, \label{cpkdv}
\end{equation} $ \alpha$ and $ \beta$ are constants.
Let \begin{equation}
q(t,x) = u(t,x)+ i v(t,x),\label{qcmplx}
\end{equation} where $ i^2=-1$. Then substituting (\ref{qcmplx}) into (\ref{cpkdv}) and separating the real and imaginary parts, we obtain \begin{align} \label{real} \Delta_1 &\equiv u_t + \alpha u u_x - \alpha v v_x + \beta u_{xxx} =0,\\
 \label{imgry} \Delta_2 & \equiv v_t + \alpha u v_x + \alpha v u_x + \beta v_{xxx}=0,
\end{align} which is a nonlinear system of coupled KdV equations. We perform Lie symmetry analysis on (\ref{real})-(\ref{imgry}), that is , we obtain Lie point symmetries,  invariant solutions and conservation laws of (\ref{real})-(\ref{imgry}).
\section{Solutions of the nonlinear coupled KdV Equations (\ref{real})-(\ref{imgry})}
\subsection{Lie point symmetries of (\ref{real})-(\ref{imgry}) \label{liepoitsymts}}
To find Lie point symmetries of the coupled KdV Equations, we use Lie's method. The infinitesimal transformations of the Lie group with parameter $\epsilon$ are 
\begin{align}
\begin{aligned}\bar{t}=t+  \xi^t(t,x,u,v)\epsilon,\,\, \bar{x}=x+  \xi^x(t,x,u,v)\epsilon,\,\,
  \bar{u}=u+  \eta^u(t,x,u,v)\epsilon,\,\, 
  \bar{v}=v+  \eta^v(t,x,u,v)\epsilon.
\label{infintli} 
\end{aligned}
\end{align}
The vector field \begin{equation}
	X= \xi^t(t,x,u,v) \frac{ \partial }{ \partial t}+ \xi^x (t,x,u,v)\frac{ \partial }{ \partial x}+ \eta^u(t,x,u,v) \frac{\partial}{ \partial u } + \eta^v(t,x,u,v) \frac{ \partial }{ \partial v},
\end{equation} is a Lie point symmetry of (\ref{real})-(\ref{imgry}) if \begin{align}\label{Eq1} \mbox{X}^{[3]}  \Delta_1\Big |_{\Delta_1=0,  \,\,\Delta_2=0} =0, \\ \label{Eq2}\mbox{X}^{[3]}  \Delta_2 \Big|_{\Delta_1=0, \,\, \Delta_2=0} =0.
\end{align} 
Expanding  (\ref{Eq1}) and (\ref{Eq2}) and splitting on derivatives of $v$ and $u$, we have an overdetermined system of ten  PDEs, namely,
\begin{align}\begin{aligned}\xi^t_u =0, \quad \xi^t_v =&0, \quad \xi^t_x =0,\quad \xi^x_u =0,\quad
 \xi^x_v =0,\quad  \xi^t_{tt}=0,\quad \xi^x_{tt}=0, \quad 3 \xi^x_x - \xi^t_t =0,\\
 3 \eta^v + 2 \xi^t_t v=&0,\quad 
  3 \alpha \eta^u + 2 \alpha \label{cp10}\xi^t_t u -3\xi^x_t =0.\end{aligned}
\end{align}
Solving the system (\ref{cp10}) in the same way we did for (\ref{1})-(\ref{8}), we obtain 
\begin{align}
\label{xit} \xi^t = & A_1+3 A_2t,\\
\label{xix} \xi^x =&A_2 x + \alpha  A_3 t + A_4,\\
\label{etau} \eta^u =&  - 2A_2 u+ A_3,\\
\label{etav} \eta^v =& -2A_2 v,
\end{align} for arbitrary constants $ A_1, A_2,A_3,A_4$.  Hence from  (\ref{xit})-(\ref{etav}),
the infinitesimal  symmetries of the coupled KdV Equations (\ref{real})-(\ref{imgry}) is a Lie algebra generated by the vector fields
\begin{align}
X_1 = &\frac{ \partial}{ \partial t}, \\
X_2 =& \frac{ \partial}{ \partial x},\\
X_3 = & \alpha t \frac{ \partial }{ \partial x}+ \frac{\partial }{ \partial u},\\
X_4 =& 3t \frac{ \partial }{ \partial t} + x\frac{ \partial}{ \partial x} - 2 u \frac{ \partial }{ \partial u}-2 v \frac{ \partial }{ \partial v}.
\end{align}
 \subsection{Commutator table}
 The set of all infinitesimal symmetries of coupled KdV equations forms a Lie algebra. By the definition of a commutator in Equation (\ref{sec2.7.2}), we get the commutation relations, which we present in Table \ref{Tab:3}.
 \begin{center}
 	
 	\begin{tabular}{|c |c|c|c|c|} 
 		\hline 
 		$[X_i,X_j]$	& $ X_1$ & $ X_2$ & $ X_3$ & $ X_4$ \\ 
 		\hline 
 		$ X_1$ & 0 & 0 & $\alpha X_2$ & $ 3X_1$ \\ 
 		\hline 
 		$ X_2$ & 0 & 0 & 0 & $ X_2$ \\ 
 		\hline 
 		$ X_3$ & -$\alpha X_2$ & 0 & 0 & -2$X_3$\\ 
 		\hline 
 		$ X_4$ & -3$X_1$ & -$ X_2$ &2$X_3$  & 0 \\ 
 		\hline 
 	\end{tabular}\captionof{table}{ \label{Tab:3} A commutator table for the Lie algebra generated by the symmetries of coupled KdV equation.}
 \end{center}
\subsection{Local Lie groups}
In this subsection, we present the corresponding group of transformations relating to each Lie point symmetry of coupled KdV system. We obtain the groups by invoking the Lie equations as defined in Theorem \ref{lieeqns}. For the system (\ref{real})-(\ref{imgry}), we have the Lie equations, for $i=1,2,3,4,$
\begin{align} \begin{aligned}
 \frac{ \mathrm{d} \bar{t}}{ \mathrm{d} \epsilon_i} = \xi^t( \bar{t},\bar{x},\bar{u},\bar{v}),\quad \bar{t}\Big |_{ \epsilon_i =0} =t, \quad
\frac{ \mathrm{d} \bar{x}}{  \mathrm{d} \epsilon_i} = \xi^x( \bar{t},\bar{x},\bar{u},\bar{v}),\quad \bar{x}\Big |_{ \epsilon_i =0} =x, \\
\frac{ \mathrm{d} \bar{u}}{  \mathrm{d} \epsilon_i} = \eta^u( \bar{t},\bar{x},\bar{u},\bar{v}),\quad \bar{u}\Big |_{ \epsilon_i =0} =u, \quad \frac{ \mathrm{d} \bar{v}}{  \mathrm{d} \epsilon_i} =\eta^v( \bar{t},\bar{x},\bar{u},\bar{v}),\quad \bar{v}\Big |_{ \epsilon_i =0} =v. \label{liq} \end{aligned}
\end{align}
 Taking for each $X_i, i =1,2,3,4, \,\,  T_{\epsilon_i}$ to be the corresponding group of transformations, we obtain from the Lie equations (\ref{liq}) the groups;
 \begin{align}
  T_{\epsilon_1}& :  \bar{t}=t+\epsilon_1, \,\, \bar{x}=x,\,\,  \bar{u}=u,\,\, \bar{v}=v,\\
  T_{\epsilon_2}& :  \bar{t}=t, \,\,  \bar{x}=x +\epsilon_2 ,\,\, \bar{u}=u, \,\, \bar{v}=v,\\ 
 T_{\epsilon_3}& :  \bar{t}=t, \,\, \bar{x}=x + \alpha\epsilon_3t,\,\, \bar{u}=u+\epsilon_3,\,\, \bar{v}=v,\\
  T_{\epsilon_4} &:  \bar{t}=t e^{3 \epsilon_4}, \,\, \bar{x}= xe^{\epsilon_4},\,\, \bar{u} = ue^{-2\epsilon_4},\,\, \bar{v}=ve^{-2\epsilon_4}. \label{lclg4}
 \end{align}
\subsection{Symmetry reductions of the coupled KdV Equations (\ref{real})-(\ref{imgry})}
We use the symmetries obtained in \textbf{Subsection \ref{liepoitsymts} Lie point symmetries of (\ref{real})-(\ref{imgry})} to perform symmetry reductions for the Coupled KdV Equations (\ref{real})-(\ref{imgry}).
\begin{enumerate} [(i)]
\item  \textbf{The  time translation symmetry $X_1 = { \partial } /{ \partial t}$. \label{stat}} \newline
Solving the characteristic equations associated to the operator $X_1$ gives the  invariants \newline$J_1 =x, \quad  J_2 = u,\quad \text{and} \quad  J_3 =v.$
 Hence, we have
 \begin{align}
 \label{gp1}	u =\varphi(x), \quad 	 v= \psi(x), 
 \end{align} for arbitrary functions $\varphi $ and  $\psi $. Substituting  the expressions  for $u$
and $v$ given by (\ref{gp1}) into the system (\ref{real})-(\ref{imgry}), we get a system of third order ordinary DEs namely, 
\begin{align} \alpha \left[  \label{krp}\varphi(x) \varphi'(x)- \psi(x) \psi'(x) \right] + \beta \varphi'''(x)=0,\\
\alpha \left( \varphi(x) \psi(x) \right)' + \beta \psi'''(x)=0. \label{spt}\end{align}
Integration of the system (\ref{krp})-(\ref{spt}) yields;
\begin{align}
\label{a} \frac{\alpha}{2}\left[  \varphi(x)^2- \psi(x)^2  \right] + \beta \varphi''(x)=C_1,\\
\label{b} \alpha \left[ \varphi(x) \psi(x)\right] + \beta \psi''(x)=C_2,
\end{align} for arbitrary constants $C_1$ and $C_2$.
If we take $C_1=C_2=0$, the system (\ref{a})-(\ref{b}) becomes   \begin{align}
\label{a1} \frac{\alpha}{2}\left[  \varphi(x)^2- \psi(x)^2  \right] + \beta \varphi''(x)=0,\\
\label{b1} \alpha \left[ \varphi(x) \psi(x)\right] + \beta \psi''(x)=0,
\end{align}

To find more solutions of the system (\ref{a1})-(\ref{b1}), we determine its Lie point symmetries.
Using the Lie's algorithm for computing point symmetries, we see that  the Lie point symmetries of (\ref{a1})-(\ref{b1}) are \begin{align} X_1^{*} = \frac{\partial }{ \partial x},\,\, X_2^{*}=x \frac{ \partial }{ \partial x}- 2 \varphi \frac{ \partial }{ \partial \varphi}-2 \psi
\frac{\partial }{ \partial \psi}.
\end{align}
Proceeding as above, we see that the symmetry $X_1^{*}$ yields the trivial solution \begin{align}
u=0,\,\, v=0. \label{tri}
\end{align}
The second symmetry $ X_2^{*}$  has the characteristic equations \begin{align}
\frac{\mathrm{d}x}{x}= \frac{\mathrm{d} \varphi}{-2 \varphi}=\frac{\mathrm{d} \psi }{-2 \psi},
\end{align} 
which provides the invariants $ J_1 = x^2 \varphi, \,\,  J_2 = x^2 \psi$.
Letting $\varphi =  {\lambda}/{x^2}, \,\,  \psi =  {\mu}/{x^2}$,  substituting the values of $\varphi$ and $\psi$ into (\ref{a1})-(\ref{b1}) and solving the resulting equations yield: 
\begin{enumerate}[(a)]
	
\item \textbf{Case one.} Taking 
 $ \mu =0$ gives $\lambda =0$ or $ \lambda = -12\beta/ \alpha$.\newline
 The case $ \lambda = 0,\,\, \mu=0$ also gives the trivial solution (\ref{tri}).
One can easily see that  $ \lambda = -12\beta/\alpha,\,\, \mu=0$ gives $\varphi=- {12 \beta}/(\alpha x^2)$, 
$\psi=0$ which  is a solution of the system (\ref{a1})-(\ref{b1}).
Hence
\begin{align}
u_1(t,x)=-\frac{12 \beta}{ \alpha x^2}, \,\, v_1(t,x)=0,
\end{align} 
is a solution of the coupled KdV system (\ref{real})-(\ref{imgry}).

\item \textbf{Case two.} Taking
$\lambda=- {6\beta}/{\alpha}, \,\, \mu = \pm   {6 \beta i}/{\alpha}$ with $i^2=-1$.  Consequently, 
\begin{align}
u_2(t,x) = - \frac{6  \beta }{ \alpha x^2},\,\,
v_2(t,x) =  \frac{ 6 i \beta }{ \alpha x^2} 
\end{align} 
and 
\begin{align}
u_3(t,x) =- \frac{6  \beta }{ \alpha x^2}, \,\,
v_3(t,x) =  -\frac{6 i \beta }{ \alpha x^2} 
\end{align} 
are solutions of the coupled KdV system (\ref{real})-(\ref{imgry}). Hence Lie group analysis has given us three steady-state solutions for the coupled KdV system (\ref{real})-(\ref{imgry}) under the time translation symmetry $X_1 = {\partial }/{ \partial t}$.

\end{enumerate}
\item \textbf{The space  translation symmetry $X_2 = {\partial }/{\partial x}.$}

Solving the characteristic equations  associated to $X_2$ gives the invariants $J_1=t, \quad J_2 = u \quad  \text{and} \newline \quad J_3 =v.$ Therefore, the group-invariant solution is
\begin{align}
u =\phi(t), \quad  v= h(t), \label{gp2}	
\end{align}  for arbitrary functions $h$ and $\phi$. Substitution of the solutions from (\ref{gp2}) into  (\ref{real})-(\ref{imgry}), we get a system of first order ordinary DEs, namely,
\begin{equation}\phi'(t)=0, \quad h'(t)=0,
\end{equation} which is integrated once with respect to $t$ to  yield,
\begin{align}
\phi(t)=C_1, \quad h(t)=C_2,
\end{align} for arbitrary constants $C_1$ and $C_2$. Consequently, the space translation group-invariant solution  of the system (\ref{real})-(\ref{imgry}) is
\begin{align}
u(t,x)=C_1,  \quad  v(t,x)=C_2.
\end{align} 
\item  \textbf{The Galilean boost symmetry} $X_3 =  \alpha t { \partial }/{ \partial x}+{\partial }/{ \partial u}.$

Solving the characteristic equations associated to Galilean boost gives  the invariants
\begin{align}
J_1 =t, \quad  J_2 =v,   \quad  J_3 = -u + \frac{x}{ \alpha t}, \,\, t \neq 0.\label{onyan}
\end{align} 

Thus  the invariant solution of (\ref{real})-(\ref{imgry}) is
\begin{align} u = \frac{x}{ \alpha t}-g(t),  \,\,   v = f(t), \,\, t\neq 0, \label{gpi3}
\end{align}  for arbitrary functions $f$ and $g$.
Substitution of the values of $u$ and $v$ from (\ref{gpi3}) into the System (\ref{real})-(\ref{imgry}), we get a nonlinear system of coupled  first order ordinary DEs, namely,
\begin{align} tg'(t)+g(t)=0,\\
tf'(t)+ f(t) =0,
\end{align} whose solutions are $ g(t) = {C_1}/{t}$ and $f(t)= {C_2}/{t} $ for arbitrary constants $ C_1$ and $C_2$.
Hence the Galilean boost group-invariant solution of the system (\ref{real})-(\ref{imgry}) is
 \begin{align}
 u(t,x) =\frac{x+A}{\alpha t}, \quad  v(t,x)= \frac{C_2}{t},
\end{align} where $ A=-\alpha C_1$ and $ t \neq 0. $

\item \textbf{The scaling  $X_4 =3t \frac{ \partial }{ \partial t} + x\frac{ \partial}{ \partial x} - 2 u \frac{ \partial }{ \partial u}-2 v \frac{ \partial }{ \partial v}$.
} \newline
By solving of the characteristic equations associated to this symmetry, we obtain the invariants 
\begin{align} J_1 = \frac{x^3}{t}, \quad  J_2 =ux^2,  \quad  J_3 = vx^2.
\end{align}
Generally, the group-invariant solution pair is 
\begin{align}
	u(t,x)=  \frac{f(\lambda) }{x^2},\quad  v(t,x)= \frac{g(\lambda) }{x^2}, \quad  \lambda=\frac{x^3}{t},
\end{align} where  the functions $ f$ and $g$ satisfy  the system of third order nonlinear  coupled  ordinary DEs
\begin{align}
2 \alpha (g^2-f^2)- \lambda^2 f'  +3 \alpha \lambda (ff'-gg') + \beta (-24f + 24 \lambda f'  + 27\lambda^3 f''' )=&0,\\
-4\alpha fg-  \lambda^2 g'  + 3 \alpha \lambda(fg)' + \beta (-24g + 24 \lambda g' +27 \lambda^3g''' )=&0.
\end{align} 
\item \textbf{Linear combination of time and space translations	$X_1 + c X_2$}. \newline
We consider a symmetry $X$, which is a linear combination of the time and space translations symmetries,  that is, 
$
X= \partial_t+ c \partial_x,
$ for a constant $c$.
The invariants associated to this symmetry $X$ are 
$
J_1 = x-ct, \,\, J_2 =u, \,\,   J_3 =v.
$
Hence, the invariant solution for the symmetry  $X$ is 
\begin{align} \label{wav1} 
u = f(x-ct), \,\, v=g(x-ct), 
\end{align}  
for arbitrary functions $f$ and $g$. Substitution of $u$ and $v$ from  (\ref{wav1}) into the system (\ref{real})-(\ref{imgry}) yields a system of nonlinear third order ordinary DEs, namely \begin{align}
-cf' (\xi)+ \alpha \left\lbrace f (\xi) f' (\xi)- g (\xi) g' (\xi) \right\rbrace  + \beta f''' (\xi)=0,\\
-cg' (\xi)+ \alpha ( f (\xi)g (\xi))'+ \beta g'''(\xi)=0,
\end{align} 
which on integrating once with respect to $\xi$ yields
\begin{align}
 \label{tr1}
 -c f + \frac{1}{2} \alpha (f^2-g^2) + \beta f''+C_1=0, \\
 \label{tr2} -cg+ \alpha f g + \beta g''+ C_2 =0,
\end{align} 
for arbitrary constants $ C_1$ and $C_2$.
\begin{rem} If we take the constants $ C_1=C_2=0$, then when the wave velocity $c=0$, we can recover the stationary solutions given in (\ref{stat}).
\end{rem}
\begin{rem} Traveling wave solutions of the system (\ref{real})-(\ref{imgry}) must satisfy the system (\ref{tr1})-(\ref{tr2}).	
\end{rem}
\end{enumerate}
\section{Conservation laws of the coupled KdV Equations (\ref{real})-(\ref{imgry})}
Computation of conservation laws for the coupled KdV Equations (\ref{real})-(\ref{imgry}) is done using two methods; the method of  multipliers  and a theorem due to Ibragimov.
\subsection{Conservation laws of (\ref{real})-(\ref{imgry}) using the multipliers}
We seek local conservation law multipliers for the system (\ref{real})-(\ref{imgry}), whose determining equations are 
\begin{align}
\label{vau}\frac{\delta}{ \delta u} \left[ \Lambda^1\Delta_1 + \Lambda^2 \Delta_2 \right] =0,\\
\label{vav}\frac{\delta}{ \delta v} \left[ \Lambda^1\Delta_1 + \Lambda^2 \Delta_2 \right] =0,
\end{align} where 
\begin{align} \frac{ \delta }{ \delta u} &= \frac{ \partial }{ \partial u}-D_t \frac{\partial }{ \partial u_t}-D_x \frac{\partial }{ \partial u_x}+ D^2_x \frac{ \partial }{ \partial u_{xx}}-D_x^3 \frac{ \partial }{ \partial u_{xxx}}+ \ldots,\\
\frac{ \delta }{ \delta v} &= \frac{ \partial }{ \partial v}-D_t \frac{\partial }{ \partial v_t}-D_x \frac{\partial }{ \partial v_x}+ D^2_x \frac{ \partial }{ \partial v_{xx}}- D_x^3\frac{ \partial }{ \partial v_{xxx}}+ \cdots,
\end{align} are the Euler-Lagrange operators and \begin{align}
\label{dt} D_t =& \frac{\partial}{\partial t} + u_t\frac{\partial}{\partial u}+ v_t\frac{\partial}{\partial v}
+ u_{tx} \frac{\partial}{\partial u_x} +  v_{tx} \frac{\partial}{\partial v_x}+ u_{tt}\frac{\partial}{\partial u_t} + v_{tt}\frac{\partial}{\partial v_t} + \cdots  , \\
\label{dx} D_x =& \frac{\partial}{\partial x} + u_x \frac{\partial}{\partial u}+ v_x \frac{\partial}{\partial v}+ u_{xx}\frac{\partial}
{\partial u_x} + v_{xx}\frac{\partial}
{\partial v_x}+ u_{tx} \frac{\partial}{\partial u_{t}} + v_{tx} \frac{\partial}{\partial v_{t}}+\cdots,
\end{align} are total derivatives operators.
We look for second order multipliers,  that is, 
\begin{align}
	\Lambda^n = \Lambda^n(t,x,u,u_x,u_{xx},v,v_x,v_{xx}), \quad n = 1,2.
\end{align}
The determining Equations  (\ref{vau})-(\ref{vav}) become
 \begin{align}
 \label{det1}\frac{\delta}{ \delta u} \left[ \Lambda^1 \{  u_t + \alpha u u_x - \alpha v v_x + \beta u_{xxx}   \}   + \Lambda^2 \{   v_t + \alpha u v_x + \alpha v u_x + \beta v_{xxx} \} \right] =0,\\
 \label{det2}\frac{\delta}{ \delta v} \left[ \Lambda^1 \{  u_t + \alpha u u_x - \alpha v v_x + \beta u_{xxx} \} + \Lambda^2 \{ v_t + \alpha u v_x + \alpha v u_x + \beta v_{xxx}  \} \right] =0.
 \end{align} 
 Expanding  (\ref{det1})-(\ref{det2})  and splitting on derivatives of $u$ and $v$ yields an overdetermined system of 22 PDEs, namely \begin{align} \begin{aligned} &\Lambda_{xx}^1=0, \,\,  \Lambda_{xx}^2=0, \,\, \Lambda_{vx}^1=0, \,\, \Lambda_{vx}^2=0, \,\,
 \Lambda_{xv_{xx}}^1=0, \,\,  \Lambda_{xv_{xx}}^2=0, \,\,  \beta \Lambda_{vv}^1-\alpha \Lambda_{v_{xx}}^2=0,
 \\
 & \beta \Lambda_{vv}^2+ \alpha \Lambda_{vv_{xx}}^1=0, \,\,
 \Lambda_{vv_{xx}}^1=0, \quad \Lambda^2_{vv_{xx}}=0, \,\, \Lambda_{v_{xx}v_{xx}}^1=0, \,\,
 \Lambda_{v_{xx}v_{xx}}^2=0, \,\,
 \Lambda^{1}_u+\Lambda^{2}_v=0, 
 \\ 
 &   \Lambda_{t}^1+ \alpha \left(\Lambda_{x}^2v +\Lambda_{x}^1u\right)=0,\,\,
 \Lambda_{t}^2+ \alpha \left( \Lambda_{x}^2u- \Lambda_{x}^1v\right)=0,  \,\,  \Lambda^{2}_u- \Lambda^{1}_v=0, \,\,\Lambda^{1}_{u_x}=0,
 \\ 
 &   
 \quad \Lambda^{2}_{u_x}=0, \,\, \Lambda^{1}_{u_{xx}}+ \Lambda^{2}_{v_{xx}}=0,  
 \Lambda^{2}_{u_{xx}}-\Lambda^{1}_{v_{xx}}=0, \,\, \Lambda^{2}_{v_{x}} =0 ,\,\, \Lambda^{1}_{v_{x}}=0 .
 \label{od1}
 \end{aligned}\end{align}
Calculations reveal the solution of the system (\ref{od1})
as  \begin{align} \begin{aligned}
\Lambda^1 =& \frac{\alpha }{ 2 \beta} \left( c_3 \{ u^2-v^2\} + 2c_4uv \right) 
+ (c_2t+c_5)u + (c_1t+c_6)v+ c_3 u_{xx}\\& + c_4v_{xx}  +c_7  -\frac{1}{\alpha} c_{2} x,\\
\Lambda^{2}=& \frac{\alpha }{2 \beta } \left(c_4 \{u^2-v^2\}  -2 c_3uv+  \right) +(c_1t + c_6)u-(c_2t +c_5)v +c_4 u_{xx} \\&-c_3 v_{xx}+c_8   - \frac{1}{ \alpha}c_1  x,\end{aligned} \end{align} for arbitrary constants $c_1,\ldots,c_8$.
\begin{rem}
Essentially, the nonlinear coupled system of KdV Equations (\ref{real})-(\ref{imgry}) have eight sets of local conservation law multipliers.
\end{rem}
\begin{rem} Recall that the multipliers must satisfy the property defined in (\ref{sec2.9.13}). For the nonlinear coupled KdV system  (\ref{real})-(\ref{imgry}), we have
	\begin{align}
	\label{oda}		\Lambda^1  \Delta_1+ \Lambda^2 \Delta_2 = D_t T^t + D_x T^x,
	\end{align} where 
	$T^t=T^t(t,x,u,v)$	and $T^x=T^x(t,x,u,v,u_x,v_x,u_{xx},v_{xx}).$	
\end{rem}
Solving (\ref{oda}) in the same way we did in  (\ref{consl}) for the KdV Equation (\ref{kd}), we obtain conserved vectors  corresponding to each set of multipliers as shown below.
\begin{enumerate}[(i)]
	\item The multiplier $ \left( \Lambda_{1}^1, \Lambda_{1}^2\right)= \left( tv,tu-\frac{x}{ \alpha}\right)$ has the  conserved vectors 
\begin{align}
T_1^t=& tuv-\frac{xv}{\alpha},\\
T_1^x=& \beta \left[ t\{vu_{xx} +uv_{xx} -v_x u_x \}   + \frac{1}{\alpha}\{v_x-xv_{xx}\} \right]+ \alpha \left[ t \left( u^2v-\frac{v^3}{3}\right) \right] -xuv.
\end{align}

\item The multiplier $ \left( \Lambda_{2}^1, \Lambda_{2}^2\right)= \left( tu-\frac{x}{ \alpha},-tv\right)$ has the  conserved vectors 
\begin{align} \begin{aligned}
T_2^t=& \frac{ t}{2} \{u^2-v^2  \} -\frac{xu}{ \alpha},\\
T_2^x=& \beta\left[t \left( uu_{xx}-vv_{xx} + \frac{ 1}{2} \{v_x^2-u_x^2 \} \right) +  \frac{1}{\alpha}\{u_x-xu_{xx} \}   \right] + \alpha t \left[ \frac{u^3}{3} -uv^2 \right]\\ & + \frac{x}{2} \{v^2-u^2 \}.\end{aligned}
\end{align}

\item  The multiplier $ \left( \Lambda_{3}^1, \Lambda_{3}^2\right)= \left(  \frac{ \alpha }{2 \beta} \{u^2- v^2 \}+ u_{xx},  -\{ \frac{ \alpha uv}{\beta} +v_{xx}\} \right)$ has the  conserved vectors 
\begin{align}
T_3^t=& \frac{\alpha }{2 \beta}
\left( \frac{u^3}{3}-uv^2\right),\\
T_3^x=& \frac{\alpha}{2} \left[ (u^2-v^2)u_{xx} -v^2v_{xx} \right] - \alpha u v v_{xx}+ \frac{\beta}{2} \left[ u_{xx}^2-v_{xx}^2\right]+ u_tu_x-v_tv_x \\& + \frac{ \alpha^2}{4 \beta} \left[ \frac{1}{2} \{ u^4+v^4 \}-3u^2v^2 \right].
\end{align}
\item  The multiplier $ \left( \Lambda_{4}^1, \Lambda_{4}^2\right)= \left( \{ \frac{ \alpha uv}{\beta} +v_{xx}\} ,  \frac{ \alpha [u^2- v^2 ] }{2 \beta}+ u_{xx}\right)$ has the  conserved vectors 
\begin{align}
T_4^t=& \frac{\alpha }{2 \beta}
\left(u^2v - \frac{v^3}{3}\right),\\
T_4^x=& \frac{ \alpha^2}{ 2 \beta} \left[  ( u^3v-uv^3)\right]+ v_tu_x + u_tv_x+ \frac{\alpha}{2} (u^2-v^2)v_{xx} + \{\alpha uv + \beta v_{xx} \}u_{xx}.
\end{align}

\item  The multiplier $ \left( \Lambda_{5}^1, \Lambda_{5}^2\right)= \left(u , -v\right)$ has the  conserved vectors 
\begin{align}
 T^t_5= \frac{1}{2}\{ u^2-v^2\},\quad
 T^x_5=\beta \left( u u_{xx} -v v_{xx} + \frac{ v_x^2-u_x^2}{2} \right)  +  \alpha \left(\frac{u^3}{3} - uv^2 \right).
\end{align}

\item  The multiplier $ \left( \Lambda_{6}^1, \Lambda_{6}^2\right)= \left(v, u\right)$ has the  conserved vectors 
\begin{align}
 T^t_6= uv,\quad
T^x_6= \beta \left( v u_{xx} +u v_{xx} -u_x v_x \right) + \alpha \left( u^2 v-  \frac{  v^3}{3} \right).
\end{align}
 \item The multiplier  $  \left( \Lambda_{7}^1, \Lambda_{7}^2\right)= \left( 1,0\right)$
 has the  conserved vectors 
 \begin{align}
 T_7^t=u,\quad
 T_7^x= \frac{\alpha}{2} \{ u^2-v^2\}+ \beta u_{xx}.
 \end{align}
\item The multiplier has  $ \left( \Lambda_{8}^1, \Lambda_{8}^2\right)= \left( 0,1\right) $ the conserved vectors 
\begin{align}
T_8^t=v,\quad
T_8^x=\alpha uv+ \beta v_{xx}.
\end{align}
\end{enumerate}
\begin{rem}
It can be verified that \begin{align}
D_t T^t_i + D_xT^x_i \Big |_{\Delta_1=0, \,\,\Delta_2=0}=0, \label{bao}
\end{align} for $i =1,\ldots,8$.
\end{rem}
\begin{rem}
The expressions in (\ref{bao}) are  eight conservation laws for the coupled KdV system (\ref{real})-(\ref{imgry}).
\end{rem}
\begin{rem}
The presence of multipliers $\left( \Lambda_{7}^1, \Lambda_{7}^2\right)= \left( 1,0\right)$ and $\left( \Lambda_{8}^1, \Lambda_{8}^2\right)= \left( 0,1\right)$ manifest that the coupled KdV equations are themselves conservation laws.
\end{rem}
\subsection{ Conservation laws of (\ref{real})-(\ref{imgry}) using Ibragimov's theorem}
In this section, we derive conserved vectors for coupled KdV equations (\ref{real})-(\ref{imgry}) by a new theorem due to Ibragimov. The concepts used here are from \textbf{Section \ref{ibra} Ibragimov's conservation theorem } in Chapter \ref{ch2}.
The adjoint equations for the nonlinear system coupled KdV Equations (\ref{real})-(\ref{imgry}) are \begin{align} \label{ad1}\Delta_1^{*}& \equiv  f_t+ \alpha \ uf_x+\alpha vg_x+ \beta f_{xxx}=0,\\
\label{ad2}\Delta_2^{*}& \equiv  g_t-\alpha  vf_x+ \alpha ug_x+ \beta g_{xxx}=0.
\end{align} The formal Lagrangian $\mathcal{L}$ for the nonlinear coupled system of the KdV Equations (\ref{real})-(\ref{imgry}) and its adjoint Equations (\ref{ad1})-(\ref{ad2}) is given by \begin{equation} \mathcal{L}= f \{u_t + \alpha u u_x- \alpha v v_x + \beta u_{xxx} \}  + g \{ v_t + \alpha  u v_x + \alpha v u_x + \beta v_{xxx}\},
\end{equation} where $f$ and $g$ are new variables.
We shall use the Lie point symmetries of the system (\ref{real})-(\ref{imgry}) ,namely \begin{align}
X_1 = \partial_t, \quad 
X_2 = \partial_x, \quad 
X_3 = \alpha t \partial_x+  \partial_u,\quad 
X_4 =3t \partial_t + x \partial_x - 2 u  \partial_u-2 v  \partial_v,
\end{align} to derive conserved vectors corresponding to each symmetry below.
\begin{enumerate} [\text{Case} (i)]
\item The symmetry $ X_1 = \partial_t$ yields Lie characteristic functions given by  $	W_1^1 =-u_t$ and \newline $W_1^2 =-v_t.$ Hence the  associated conserved vector is given by \begin{align} \begin{aligned}
T_1^t=& \alpha \left[f \{ u u_x- v v_x \} + g \{v u_x + u v_x \}  \right]    + \beta  \{ f u_{xxx}+ gv_{xxx} \} ,\\
T_1^x=&  \alpha \left[ f \{-  uu_t+vv_t \} - g\{ v u_t+uv_t\}\right]  \\& +  \beta \{ f_x u_{tx} + g_x v_{tx} - u_t f_{xx} -v_t g_{xx} - fu_{txx} - gv_{txx}\}. \end{aligned}\end{align}
\item The symmetry $ X_2 = \partial_x$ yields Lie characteristic functions  $
W_2^1 =-u_x, \,\, \text{and}$ $  W_2^2 =-v_x$. Therefore the associated conserved vector is
\begin{align}
T_2^t&=- u_x f - v_xg,\\
 T_2^x&= fu_t + gv_t + \beta\{  -u_xf_{xx} -v_x g_{xx} + f_{x}u_{xx} + g_{x} v_{xx}\}. 
\end{align}
\item The symmetry $ X_3=  \alpha t \partial_x+  \partial_u$ yields Lie characteristic functions given by \newline$W_3^1 =1-\alpha t u_x, \,\, \text{and} \,\, W_3^2 =-\alpha tv_x.$
Hence the associated  conserved vector is given by \begin{align} \begin{aligned}
T_3^t= &f- \alpha t \{  u_x f +  v_x g\},\\
T_3^x=& \alpha \left[ fu+gv+ t\{ u_tf+v_tg\} + \beta t\{ \frac{f_{xx}}{ \alpha t} -u_xf_{xx} -v_x g_{xx} + f_x u_{xx} + g_x v_{xx}    \}\right]. \end{aligned}
\end{align}

\item The symmetry $ X_4 = 3t  \partial_t + x \partial_x - 2 u \partial_u-2 v  \partial_v$ yields  the Lie characteristic functions: 
$W_4^1 =-2u -3tu_t-xu_x, \quad  W_4^2 =-2v -3tv_t-xv_x.
$ Consequently, the  corresponding conserved vector is given by
\begin{align} \begin{aligned} T_4^t&=\alpha \left[ 3t \{fuu_x-fvv_x+guv_x+gvu_x \}  \right]+ \beta \left[ 3t\{ fu_{xxx} + gv_{xxx}\}\right]\\& -2\{fu+gv \}-x\{ fu_x+gv_x\},
\\
T_4^x& = x\{fu_t+gv_t \}+ \beta\left[ 3\begin{pmatrix} f_xu_x+g_xv_x+t\{ f_xu_{tx} +g_x v_{tx} \} \end{pmatrix}\right]\\& 
-\alpha \left[2 \left(f \{u^2-v^2 \}+2guv \right)+3t\begin{pmatrix} f\{ uu_t-vv_t\}+g\{vu_t+uv_t\} \end{pmatrix}\right] \\&
-\beta \left[ x \{ u_xf_{xx}+v_xg_{xx}-f_xu_{xx}-g_xv_{xx} \} +2 \{ uf_{xx} + vg_{xx}\}\right]\\&
-\beta \left[ 3t \{ f_{xx}u_t+g_{xx}v_t+ fu_{txx} + gv_{txx}\} + 4 \{ fu_{xx} + gv_{xx}\}\right].
\end{aligned}
\end{align} 
\end{enumerate}
\begin{rem}
The appearance of arbitrary functions $f(t,x)$ and $g(t,x)$ in the conserved vectors proves the existence of infinite conservation laws for coupled KdV system obtained by Ibagimov's method.
\end{rem}
\section{Concluding remarks}
In this chapter, we have studied the nonlinear coupled system of  KdV equations by use of Lie group analysis. Just like for the KdV 
Equation (\ref{kd}), we have symmetries that represent time and space translations, Galilean boost, and scaling for the nonlinear coupled KdV equation. For each symmetry, we performed a  reduction of the nonlinear coupled system.
All the group-invariant solutions describe the various states of the system.  Last but not least, an infinite number of conservation laws were derived for the system by the multiplier approach and Ibragimov's conservation theorem.

Three of the conservation laws, just like those for the KdV equation in Chapter \ref{ch3} manifest that mass, momentum, and energy are invariant quantities in the evolution of the coupled KdV system. In fact, only some of the first laws have a physical interpretation. Higher-order laws aid in understanding the qualitative properties of solutions. These conservation laws are very important in explaining the integrability of a system and the effectiveness of numerical methods used in approximating solutions. They can also be used to get exact solutions to the nonlinear coupled KdV system, which is a subject of future work. % Conclusion is usually a chapter itself. 
\chapter{Conclusion \label{ch5}}
This dissertation aimed to research on Lie group analysis of nonlinear partial differential equations (NLPDEs). It must be said that NLPDEs are very important in describing various phenomena and systems in real life. Our main concern in this study was the nonlinear coupled system of Korteweg-de Vries (KdV) equations that describe the dynamics of solitons. However, this powerful approach that uses invariant properties of an NLPDE to get exact solutions and conservation laws can be used for any NLPDE.

The first step was to develop mathematical concepts of Lie group analysis and conservation laws in Chapter \ref{ch2}. We developed pertinent notions of local Lie groups, prolongations, symmetry groups and invariants. Also introduced is the concept of conservation laws and two methods to compute them, that is, the method of multipliers and Ibragimov's conservation theorem.

Thereafter, we studied the KdV equation as an illustrative example in Chapter \ref{ch3}. This involved finding its space and time translations, Galilean boost, and scaling symmetries. We then performed symmetry reductions that resulted in a total of four group-invariant solutions. A linear combination of time and space translations also helped in finding the one-soliton solution for the KdV equation. Finally, we found conservation laws for mass, momentum, and energy by using the multipliers method.


Our third step was to investigate a system of coupled KdV equations in Chapter \ref{ch4}. Likewise to the KdV equation studied in Chapter \ref{ch3}, a four-dimensional Lie algebra of symmetries was found for the nonlinear coupled system KdV equations. This was also spanned by  space and time translations, Galilean boost and scaling symmetries where the scaling symmetry acts on four variables. Associated to each symmetry, we obtained symmetry reductions that gave six nontrivial solutions for the coupled system.  Lastly, we constructed infinite conservation laws of a nonlinear coupled KdV system by using multipliers and a theorem proposed by Nail Ibragimov. Three of these laws show that mass, momentum and energy are conserved quantities in the evolution of a nonlinear coupled KdV system.

The above results show a very interesting property of the KdV equation. Most important to note is that the infinite number of conservation laws for the coupled system show that the KdV equation is completely integrable, meaning that the behavior of the system can be determined by initial conditions and can be integrated from the prescribed initial conditions. Indeed, the KdV equation gives rise to multiple-soliton solutions thus emphasizing the importance of the KdV equation in the theory of integrable systems. The beautiful KdV equation is ubiquitous, having applications in various settings.
The equation has been used to describe the dynamics of solitons, ion-acoustic waves in a plasma, shallow-water waves  and 
nonlinear perturbations along internal surfaces between layers of different densities in stratified fluids, for example propagation of solitons of long internal waves in oceans. The KdV equation also models shock wave formation, turbulence, boundary layer behavior, and mass transport.


\section*{Further work}
 For future work, we are interested in applying the conservation laws to construct exact solutions to the nonlinear coupled system of Korteweg-de Vries equations. % You may have more chapters. (Use e.g. git add FILE)
% This is where we stop counting pages.
% Acknowledgements and References are not counted.
%-----------------------------------------------------------------------------
% See the acknowledgement.tex file and follow the instructions there.
\chapter*{Acknowledgements}
% Don't change anything above this.
% We do not number this or add it to the contents!
% Overly long acknowledgements are not professional.

I want to acknowledge AIMS and its funders for supporting this work, as well as my supervisor, Prof Chaudry Masood Khalique from North-West University.



%-----------------------------------------------------------------------------
% Note the errata page is not for now, it is for use during the examination.
% Not that you're going to have any errata.
%-----------------------------------------------------------------------------
% THE BIBLIOGRAPHY 
% Bibliography styles define how the bibliography is 
% listed and formatted. This is part of the AIMS house
% style and is only changed under exceptional circumstances
\renewcommand{\bibname}{References}
\nocite{*}
\bibliographystyle{myabbrvnat}
\bibliography{references}
\addcontentsline{toc}{chapter}{References}
%-----------------------------------------------------------------------------

\end{document}
